% =============================================================================
% BACHELORARBEIT
% Entwicklung und Evaluation einer zustandsorientierten Agenten-Architektur
% durch spezialisierte Custom-Nodes in n8n
% =============================================================================

\documentclass[12pt,a4paper,oneside]{report}

% ---------------------------------------------------------------------------
% Pakete
% ---------------------------------------------------------------------------
\usepackage[T1]{fontenc}
\usepackage[utf8]{inputenc}
\usepackage[ngerman]{babel}
\usepackage{geometry}
\usepackage{setspace}
\usepackage{graphicx}
\usepackage{listings}
\usepackage{xcolor}
\usepackage{hyperref}
\usepackage{booktabs}
\usepackage{tabularx}
\usepackage{longtable}
\usepackage{array}
\usepackage{enumitem}
\usepackage{caption}
\usepackage{subcaption}
\usepackage{amsmath}
\usepackage{amssymb}
\usepackage[protrusion=true,expansion=false]{microtype}
\usepackage{csquotes}
\usepackage{parskip}
\usepackage{fancyhdr}
\usepackage{titlesec}
\usepackage{tocloft}
\usepackage{acronym}
\usepackage{float}
\usepackage{wrapfig}

% ---------------------------------------------------------------------------
% Seitenlayout
% ---------------------------------------------------------------------------
\geometry{
  a4paper,
  left=3.0cm,
  right=2.5cm,
  top=2.5cm,
  bottom=2.5cm,
  headsep=1cm,
  headheight=14.5pt
}

\onehalfspacing

% ---------------------------------------------------------------------------
% Kapitel-Formatierung (verhindert den grossen Leerraum oben bei \chapter)
% ---------------------------------------------------------------------------
\titleformat{\chapter}[display]
  {\normalfont\huge\bfseries}
  {\chaptertitlename\ \thechapter}
  {10pt}
  {\Huge}
\titlespacing*{\chapter}{0pt}{0pt}{20pt}

% Abschnitte etwas kompakter
\titlespacing*{\section}{0pt}{16pt}{6pt}
\titlespacing*{\subsection}{0pt}{12pt}{4pt}
\titlespacing*{\subsubsection}{0pt}{10pt}{2pt}

% ---------------------------------------------------------------------------
% Kopf- und Fußzeile
% ---------------------------------------------------------------------------
\pagestyle{fancy}
\fancyhf{}
\fancyhead[L]{\nouppercase{\leftmark}}
\fancyhead[R]{\thepage}
\renewcommand{\headrulewidth}{0.4pt}

% ---------------------------------------------------------------------------
% Code-Listing-Stil
% ---------------------------------------------------------------------------
\definecolor{codegray}{rgb}{0.5,0.5,0.5}
\definecolor{codepurple}{rgb}{0.58,0,0.82}
\definecolor{backcolour}{rgb}{0.97,0.97,0.97}
\definecolor{codegreen}{rgb}{0,0.5,0}
\definecolor{navyblue}{rgb}{0.0,0.0,0.5}

\lstdefinestyle{typescript}{
  backgroundcolor=\color{backcolour},
  commentstyle=\color{codegreen},
  keywordstyle=\color{navyblue}\bfseries,
  numberstyle=\tiny\color{codegray},
  stringstyle=\color{codepurple},
  basicstyle=\ttfamily\footnotesize,
  breakatwhitespace=false,
  breaklines=true,
  captionpos=b,
  keepspaces=true,
  numbers=left,
  numbersep=5pt,
  showspaces=false,
  showstringspaces=false,
  showtabs=false,
  tabsize=2,
  frame=single,
  rulecolor=\color{gray!40},
}

\lstset{style=typescript}

% ---------------------------------------------------------------------------
% Hyperref-Einstellungen
% ---------------------------------------------------------------------------
\hypersetup{
  colorlinks=true,
  linkcolor=navyblue,
  citecolor=navyblue,
  urlcolor=navyblue,
  pdftitle={Entwicklung und Evaluation einer zustandsorientierten Agenten-Architektur durch spezialisierte Custom-Nodes in n8n},
  pdfauthor={[Name eintragen]},
}

% =============================================================================
% BEGINN DES DOKUMENTS
% =============================================================================
\begin{document}

% ---------------------------------------------------------------------------
% Titelseite
% ---------------------------------------------------------------------------
\begin{titlepage}
  \centering
  \vspace*{1cm}

  {\Large \textbf{[Universität Bremen]}}\\[0.5cm]
  {\large Fachbereich Informatik}\\[2cm]

  \rule{\textwidth}{1.5pt}\\[0.4cm]
  {\huge \textbf{Entwicklung und Evaluation einer\\[0.3cm]
  zustandsorientierten Agenten-Architektur\\[0.3cm]
  durch spezialisierte Custom-Nodes in n8n}}\\[0.4cm]
  \rule{\textwidth}{1.5pt}\\[2cm]

  {\large \textbf{Bachelorarbeit}}\\[0.5cm]
  {\normalsize im Studiengang Vollfach Informatik}\\[2cm]

  \begin{tabular}{ll}
    Vorgelegt von:    & Max Lehmann\\
    Matrikelnummer:   & 6187272\\
    Erstgutachter:    & Prof. Dr. Sebastian Maneth\\
    Zweitgutachter:   & Noch Unbekannt\\
    Abgabedatum:      & 01. Juni 2026\\
  \end{tabular}

  \vfill
  {\small \today}
\end{titlepage}

% ---------------------------------------------------------------------------
% Eidesstattliche Erklärung
% ---------------------------------------------------------------------------
\thispagestyle{empty}
\chapter*{Eidesstattliche Erklärung}
\addcontentsline{toc}{chapter}{Eidesstattliche Erklärung}

Ich erkläre hiermit an Eides statt, dass ich die vorliegende Bachelorarbeit
selbstständig und ohne unzulässige fremde Hilfe angefertigt habe. Alle verwendeten
Quellen sind vollständig angegeben; wörtlich oder sinngemäß übernommene Stellen sind
als solche kenntlich gemacht.

\vspace{2cm}
\noindent
Bremen, 08.04.2026

\vspace{1.5cm}
\noindent\rule{7cm}{0.4pt}\\
Max Lehmann

\newpage

% ---------------------------------------------------------------------------
% Kurzfassung / Abstract
% ---------------------------------------------------------------------------
\thispagestyle{empty}
\chapter*{Kurzfassung}
\addcontentsline{toc}{chapter}{Kurzfassung}

Die vorliegende Bachelorarbeit entwickelt und evaluiert eine zustandsorientierte
Agenten-Architektur für die Low-Code-Plattform n8n. Auf Basis der Design-Science-Research-Methodik
nach Hevner et al.\ (2004) werden drei spezialisierte Custom Nodes implementiert:
\textit{AgentStateNode}, \textit{VectorStoreNode} und \textit{ReActLoopNode}. Diese
realisieren gemeinsam das ReAct-Muster (Yao et al., 2022) innerhalb des visuellen
n8n-Workflow-Editors und ermöglichen so transparente, kontrollierbare und in
Unternehmenssysteme integrierbare KI-Agenten.

Im Gegensatz zu autonomen Agenten-Frameworks wie LangChain, AutoGPT oder Claude Code,
die Entscheidungsprozesse in einer Black Box verbergen, macht das entwickelte Artefakt
jeden Thought-Action-Observation-Zyklus als diskreten Workflow-Schritt sichtbar.
Die Auswertung von drei Benchmark-Szenarien zeigt, dass das Artefakt bei Transparenz
und Kontrollierbarkeit klare Vorteile gegenüber der Standard-n8n-Lösung aufweist. Diese
Vorteile gehen jedoch mit einem größeren Workflow und teils höherer Latenz einher; in
Szenarien ohne komplexe Zustandsverwaltung ist der schlankere Basis-Ansatz weiterhin
die einfachere Wahl.

\vspace{1cm}
\noindent\textbf{Schlüsselwörter:} KI-Agenten, n8n, Custom Nodes, ReAct, RAG,
zustandsorientierte Architektur, Low-Code, LLM, Design Science Research

\vspace{1.5cm}
\noindent\textbf{Abstract (English)}

This bachelor's thesis develops and evaluates a stateful agent architecture for the
low-code platform n8n. Following the Design Science Research methodology of Hevner et al.\
(2004), three specialised custom nodes are implemented: \textit{AgentStateNode},
\textit{VectorStoreNode}, and \textit{ReActLoopNode}. Together, they realise the
ReAct pattern (Yao et al., 2022) within n8n's visual workflow editor, enabling
transparent, controllable, and enterprise-integrable AI agents.

Unlike autonomous agent frameworks such as LangChain, AutoGPT, or Claude Code, which
conceal decision processes inside a black box, the developed artefact makes every
Thought-Action-Observation cycle visible as a discrete workflow step. The evaluation of
three benchmark scenarios shows that the artefact achieves higher transparency and
controllability scores, though at the cost of greater workflow complexity and, in some
scenarios, longer execution times.

\newpage

% ---------------------------------------------------------------------------
% Inhaltsverzeichnis
% ---------------------------------------------------------------------------
\tableofcontents
\newpage

% ---------------------------------------------------------------------------
% Abbildungsverzeichnis
% ---------------------------------------------------------------------------
\listoffigures
\addcontentsline{toc}{chapter}{Abbildungsverzeichnis}
\newpage

% ---------------------------------------------------------------------------
% Tabellenverzeichnis
% ---------------------------------------------------------------------------
\listoftables
\addcontentsline{toc}{chapter}{Tabellenverzeichnis}
\newpage

% ---------------------------------------------------------------------------
% Abkürzungsverzeichnis
% ---------------------------------------------------------------------------
\chapter*{Abkürzungsverzeichnis}
\addcontentsline{toc}{chapter}{Abkürzungsverzeichnis}

\begin{acronym}[HITL]
  \acro{API}{Application Programming Interface}
  \acro{CoT}{Chain-of-Thought}
  \acro{DSR}{Design Science Research}
  \acro{ERP}{Enterprise Resource Planning}
  \acro{HITL}{Human-in-the-Loop}
  \acro{JSON}{JavaScript Object Notation}
  \acro{KI}{Künstliche Intelligenz}
  \acro{LLM}{Large Language Model}
  \acro{RAG}{Retrieval-Augmented Generation}
  \acro{SDK}{Software Development Kit}
  \acro{SQL}{Structured Query Language}
  \acro{TAO}{Thought-Action-Observation}
  \acro{UUID}{Universally Unique Identifier}
\end{acronym}

\newpage

% =============================================================================
% KAPITEL 1: EINLEITUNG
% =============================================================================
\chapter{Einleitung}

\section{Motivation und Problemstellung}

Sprachmodelle (\acf{LLM}) sind in den letzten Jahren zunehmend in betriebliche
Softwareprozesse eingeflossen und werden inzwischen auch in der Unternehmens-IT
ernsthaft diskutiert. Systeme wie AutoGPT, LangChain Agents
oder der native \textit{AI Agent Node} von n8n versprechen eine weitgehend selbstständige
Aufgabenbearbeitung, indem sie \ac{LLM}s mit externen Werkzeugen -- sogenannten Tools
-- verbinden und so mehrstufige Aufgaben autonom abarbeiten können \cite{wang2023survey}.

In der Praxis offenbaren diese Ansätze jedoch erhebliche Schwächen, sobald sie in
produktive Unternehmensumgebungen überführt werden sollen:

\begin{itemize}
  \item \textbf{Fehlende Transparenz:} Autonome Agenten wie AutoGPT oder Claude Code
    treffen ihre Entscheidungen intern, ohne dass ein menschlicher Operator die einzelnen
    Reasoning-Schritte nachvollziehen kann. Audit-Anforderungen und regulatorische
    Vorgaben -- etwa im Finanz- oder Gesundheitsbereich -- sind damit schwer zu erfüllen.

  \item \textbf{Mangelnde Kontrollierbarkeit:} Sobald ein Agenten-Framework gestartet
    ist, fehlen granulare Möglichkeiten, einzelne Entscheidungen zu genehmigen oder
    abzulehnen. Dies ist insbesondere bei irreversiblen Aktionen -- wie dem Schreiben
    in eine Produktionsdatenbank oder dem Versenden von E-Mails -- problematisch.

  \item \textbf{Schlechte Prozessintegration:} Bestehende Frameworks setzen auf
    Python-Code und eigene Bibliotheken. Die Integration in bestehende IT-Landschaften
    (ERP-Systeme, SQL-Datenbanken, REST-\ac{API}s) erfordert erheblichen Entwicklungsaufwand
    und ist nicht ohne Weiteres für Nicht-Programmierer zugänglich.

  \item \textbf{Zustandslosigkeit in Workflow-Engines:} Standard-Workflow-Orchestratoren
    wie n8n behandeln jeden Ausführungsschritt isoliert. Es existiert kein nativer
    Mechanismus, um den Fortschritt eines mehrstufigen Agenten über mehrere
    Node-Ausführungen hinweg zu speichern und bei Fehlern wiederaufzunehmen.
\end{itemize}

Daraus ergibt sich die Fragestellung dieser Arbeit: Lässt sich ein LLM-gesteuerter
Agent innerhalb einer bestehenden Low-Code-Plattform betreiben, der seinen Zustand
über mehrere Schritte hinweg behält, jeden Entscheidungsschritt nachvollziehbar macht
und dabei die vorhandenen Integrationsmöglichkeiten der Plattform nutzen kann?

\section{Zielsetzung}

Ziel dieser Arbeit ist die Entwicklung und Evaluation eines \textit{n8n Agentic Pack}
-- eines Pakets aus drei spezialisierten Custom Nodes, das eine zustandsorientierte
Agenten-Architektur innerhalb des n8n-Workflow-Editors realisiert. Die drei Nodes
bilden ein koordiniertes System:

\begin{enumerate}
  \item \textbf{AgentStateNode:} Implementiert ein persistentes Zustandsmodell für
    den laufenden Agenten-Prozess, basierend auf dem Memory-Stream-Konzept von Park
    et al.\ \cite{park2023generative}.

  \item \textbf{VectorStoreNode:} Stellt eine native ChromaDB-Integration für die
    Einbettung, Speicherung und semantische Suche von Dokumenten bereit und realisiert
    damit die Retriever-Komponente einer \ac{RAG}-Architektur nach Lewis et al.\
    \cite{lewis2020retrieval}.

  \item \textbf{ReActLoopNode:} Orchestriert Thought\,$\to$\,Action\,$\to$\,Observation-Zyklen
    nach dem ReAct-Muster (Yao et al., 2022) und macht jeden Schritt als diskretes
    Workflow-Element im n8n-Canvas sichtbar \cite{yao2022react}.
\end{enumerate}

\section{Forschungsfrage}

\begin{displayquote}
\textit{Wie können spezialisierte Custom Nodes eine effiziente, zustandsorientierte
Agenten-Architektur in n8n ermöglichen, und welche messbaren Vorteile ergeben sich
hinsichtlich Transparenz und Prozessintegration gegenüber Standardimplementierungen
und autonomen Agenten-Werkzeugen?}
\end{displayquote}

Die Forschungsfrage lässt sich in drei Teilfragen untergliedern:

\begin{description}
  \item[TF1 -- Machbarkeit:] Lässt sich das ReAct-Muster vollständig und korrekt
    innerhalb des n8n-Custom-Node-Mechanismus implementieren?
  \item[TF2 -- Transparenz:] Erzeugt die Artefakt-Implementierung eine nachvollziehbar
    höhere Transparenz der Agenten-Entscheidungen gegenüber der Baseline?
  \item[TF3 -- Integration:] Reduziert die native SDK-Integration des VectorStoreNode
    den Implementierungsaufwand gegenüber der HTTP-Workaround-Baseline messbar?
\end{description}

\section{Wissenschaftliche Einordnung}

Die Arbeit folgt der \acf{DSR}-Methodik nach Hevner et al.\ \cite{hevner2004design},
die IT-Artefakte als primäres Forschungsergebnis definiert und deren rigorose
Evaluation gegenüber einer Baseline einfordert. Hevner et al.\ formulieren sieben
Leitlinien für DSR-Forschung; die vorliegende Arbeit adressiert insbesondere:

\begin{itemize}
  \item \textbf{Leitlinie 1 -- Design als Artefakt:} Das n8n-Agentic-Pack ist das
    primäre Forschungsergebnis.
  \item \textbf{Leitlinie 2 -- Problementlevanz:} Transparenz und Kontrollierbarkeit
    von KI-Agenten sind aktuelle Herausforderungen in der Praxis.
  \item \textbf{Leitlinie 3 -- Design-Evaluation:} Drei Benchmark-Szenarien vergleichen
    Artefakt, Baseline und Referenz-Implementierung.
  \item \textbf{Leitlinie 6 -- Design als Suchprozess:} Die iterative Entwicklung
    der drei Nodes folgt dem DSR-Zyklus aus Aufbau, Evaluation und Verfeinerung.
\end{itemize}

\section{Aufbau der Arbeit}

\begin{table}[H]
  \centering
  \caption{Struktur der Bachelorarbeit}
  \label{tab:aufbau}
  \begin{tabularx}{\textwidth}{clX}
    \toprule
    \textbf{Kapitel} & \textbf{Titel} & \textbf{Inhalt} \\
    \midrule
    2 & Theoretische Grundlagen  & LLMs, ReAct, RAG, Memory, n8n-Abgrenzung,
                                   Agenten-Vergleich, \ac{DSR} \\
    3 & AgentStateNode           & Design, Implementierung, Tests, Designalternativen \\
    4 & VectorStoreNode          & Design, Implementierung, Tests, Plattformvergleich \\
    5 & ReActLoopNode            & Design, Implementierung, Tests, Architekturvergleich \\
    6 & Evaluation               & Szenarien, Metriken, Ergebnisse, Hypothesentest \\
    7 & Diskussion und Fazit     & Beantwortung der Forschungsfrage, Limitationen, Ausblick \\
    \bottomrule
  \end{tabularx}
\end{table}

% =============================================================================
% KAPITEL 2: THEORETISCHE GRUNDLAGEN
% =============================================================================
\chapter{Theoretische Grundlagen}

\section{Große Sprachmodelle als kognitive Agenten-Basis}

Große Sprachmodelle (englisch: \textit{Large Language Models}, \ac{LLM}s) sind
auf Basis der Transformer-Architektur (Vaswani et al., 2017) trainierte neuronale
Netze, die auf Milliarden von Textparametern trainiert wurden. Ihr zentraler
Mechanismus -- der \textit{Self-Attention}-Algorithmus -- ermöglicht die kontextuelle
Modellierung beliebig langer Eingabesequenzen und bildet damit die technische
Grundlage für komplexes Reasoning über mehrere Schritte hinweg \cite{openai2023gpt4}.

Für den Einsatz als Agenten-Fundament sind drei Eigenschaften moderner \ac{LLM}s
besonders relevant:

\begin{enumerate}
  \item \textbf{Instruction Following:} Modelle wie GPT-4o oder Claude 3.5 Sonnet
    können strukturierte Anweisungen -- etwa Prompt-Templates mit festen Ausgabeformaten
    -- zuverlässig befolgen. Dies ist Voraussetzung für das deterministische Parsing der
    ReAct-Ausgaben.

  \item \textbf{In-Context Learning:} \ac{LLM}s können aus wenigen Beispielen im
    Prompt (Few-Shot) verallgemeinern, ohne erneut trainiert zu werden. Das ReAct-Muster
    nutzt dies, indem vergangene Thought-Action-Observation-Tripel als Kontext mitgegeben
    werden.

  \item \textbf{Tool Use / Function Calling:} Neuere \ac{LLM}s unterstützen natives
    Function Calling via \ac{API}, bei dem das Modell strukturierte JSON-Ausgaben für
    Werkzeugaufrufe produziert. Die vorliegende Arbeit nutzt stattdessen ein
    textbasiertes ReAct-Format, das modellunabhängig ist.
\end{enumerate}

Turing (1950) formulierte bereits die philosophische Grundfrage, ob Maschinen denken
können \cite{turing1950computing}. Moderne \ac{LLM}s realisieren diese Fähigkeit in
einer praktisch nutzbaren Form: Sie ermöglichen durch die Kombination von Reasoning
und Werkzeugnutzung die Lösung komplexer, mehrstufiger Aufgaben.

\section{Das ReAct-Muster}
\label{sec:react}

\subsection{Grundkonzept und Motivation}

Das ReAct-Framework wurde von Yao et al.\ (2022) als Synthese zweier bis dahin
getrennter Forschungsrichtungen entwickelt: dem Chain-of-Thought-Prompting (Wei et al.,
2022) einerseits und der LLM-gesteuerten Aktionsausführung andererseits \cite{yao2022react}.

\subsubsection{Chain-of-Thought als Vorgänger}

Wei et al.\ (2022) zeigten, dass das explizite Formulieren von Zwischenschritten
(\textit{Chain-of-Thought}, \ac{CoT}) die Leistungsfähigkeit von \ac{LLM}s bei
komplexen Aufgaben erheblich verbessert \cite{wei2022chain}. Ein \ac{CoT}-Prompt
enthält beispielsweise:

\begin{lstlisting}[caption={Beispiel eines Chain-of-Thought-Prompts}, label=lst:cot]
Frage: Marie hat 3 Mal so viele Aepfel wie Peter.
       Peter hat 5 Aepfel. Wie viele hat Marie?
Gedankengang: Peter hat 5 Aepfel. Marie hat 3 * 5 = 15 Aepfel.
Antwort: 15
\end{lstlisting}

Der entscheidende Nachteil von \ac{CoT} liegt darin, dass das Modell ausschließlich
auf sein internes, in den Gewichten gespeichertes Wissen zurückgreift. Bei
wissensintensiven Aufgaben -- wie dem Abrufen aktueller Datenbank-Inhalte oder dem
Zugriff auf unternehmensinterne Dokumente -- führt dies systematisch zu Halluzinationen
oder veralteten Informationen.

\subsubsection{Das ReAct-Trajekt-Format}

ReAct löst dieses Problem, indem es Reasoning-Schritte (\textit{Thoughts}) mit
konkreten Werkzeugaufrufen (\textit{Actions}) und externem Feedback (\textit{Observations})
verknüpft. Yao et al.\ (2022) definieren die ReAct-Trajektorie als Sequenz von Tripeln
$(t_i, a_i, o_i)$, wobei gilt:

\begin{itemize}
  \item $t_i$ -- Thought: explizite Reasoning-Spur des Modells vor der Entscheidung
  \item $a_i$ -- Action: gewähltes Werkzeug und seine Eingabeparameter
  \item $o_i$ -- Observation: Ergebnis der Werkzeugausführung aus der Umgebung
\end{itemize}

\begin{lstlisting}[caption={ReAct-Trajekt für das SC-01-Szenario (eigene Darstellung)},
  label=lst:react-example]
Thought: Ich muss zunaechst die ausstehenden Bestellungen aus der Datenbank
         abrufen. Dazu verwende ich das sql_query-Tool.
Action: sql_query
Action Input: {"query": "SELECT * FROM orders WHERE status = 'pending'"}

Observation: [{"order_id": 1001, "customer": "Acme GmbH", "total": 1200.50},
              {"order_id": 1002, "customer": "Beta AG",   "total":  870.00}]

Thought: Ich habe 2 Bestellungen abgerufen. Nun muss ich sie in das
         ERP-Format transformieren und uebertragen.
Action: erp_push
Action Input: {"payload": [{"order_id": 1001, ...}, {"order_id": 1002, ...}]}

Observation: {"status": "ok", "recordsImported": 2}

Thought: Alle Bestellungen wurden erfolgreich importiert.
Final Answer: Beide ausstehenden Bestellungen wurden erfolgreich ins ERP
              uebertragen (Acme GmbH: 1.200,50 EUR; Beta AG: 870,00 EUR).
\end{lstlisting}

\subsection{Empirische Belege}

Yao et al.\ (2022) evaluieren ReAct auf den Benchmarks HotpotQA (mehrstufige
Frage-Antwort), Fever (Faktenverifikation) und ALFWorld (interaktive Textspiele).
In allen drei Domänen übertrifft ReAct reine \ac{CoT}-Ansätze signifikant, weil
externe Informationsquellen Halluzinationen verhindern \cite{yao2022react}. Diese
empirischen Belege motivieren die Wahl des ReAct-Musters als Grundlage des Artefakts.

\section{Agentic Memory und Zustandspersistenz}
\label{sec:memory}

\subsection{Das Memory-Stream-Modell}

Park et al.\ (2023) präsentierten generative Agenten -- simulierte menschliche
Charaktere in einer virtuellen Umgebung --, die glaubwürdiges, kohärentes Verhalten
über lange Zeiträume hinweg zeigen \cite{park2023generative}. Das Kernkonzept ist der
\textit{Memory Stream}: eine chronologisch geordnete, append-only-Liste aller erlebten
Ereignisse.

Der Memory Stream unterscheidet sich grundlegend von einem einfachen Freitext-Kontext:

\begin{table}[H]
  \centering
  \caption{Vergleich: Freitext-Kontext vs. Memory Stream nach Park et al. (2023)}
  \label{tab:memory-comparison}
  \begin{tabularx}{\textwidth}{lXX}
    \toprule
    \textbf{Eigenschaft} & \textbf{Freitext-Kontext} & \textbf{Memory Stream} \\
    \midrule
    Struktur       & Unstrukturierter Text       & Typisierte Einträge (Thought, Action, Observation) \\
    Auditierbarkeit & Schwer maschinell auswertbar & Direkt maschinenlesbar \\
    Zeitstempel    & Optional / informell          & ISO-8601, verpflichtend \\
    Filterung      & Nur manuell                   & Selektiver Abruf nach Relevanz \\
    Fehlertoleranz & Gesamter Kontext verloren     & Persistenz über Fehler hinweg \\
    \bottomrule
  \end{tabularx}
\end{table}

\subsection{Taxonomy agentischer Gedächtnistypen}

Wang et al.\ (2023) liefern eine umfassende Taxonomy von Agenten-Architekturen und
unterscheiden vier Gedächtnistypen: \textit{Sensory Memory} (kurzfristige Wahrnehmung),
\textit{Short-Term Memory} (Arbeitsgedächtnis, begrenzt), \textit{Long-Term Memory}
(extern persistiert) und \textit{Episodic Memory} (erlebte Ereignisse)
\cite{wang2023survey}. Der \textit{AgentStateNode} dieser Arbeit implementiert den
\textit{Episodic-Memory}-Typ: jede \ac{TAO}-Episode wird als timestampierter Eintrag
im History-Array gespeichert.

\section{Retrieval-Augmented Generation (RAG)}
\label{sec:rag}

\subsection{Architektur und Motivation}

Lewis et al.\ (2020) führten \ac{RAG} ein, um die faktische Halluzination von \ac{LLM}s
durch die Einbindung externer Dokumentenquellen zur Inferenzzeit zu reduzieren
\cite{lewis2020retrieval}. Die Architektur gliedert sich in zwei Phasen:

\paragraph{Indexierungsphase (offline):}
Dokumente werden in Dense-Vektoren transformiert und in einem Vektorspeicher abgelegt.
Als Embedding-Modell wird in dieser Arbeit \texttt{text-embedding-3-small} von OpenAI
verwendet, das 1.536-dimensionale Vektoren produziert.

\paragraph{Retrieval-Phase (online):}
Eine Anfrage wird ebenfalls eingebettet; der Vektorspeicher gibt die $k$ ähnlichsten
Dokumente nach dem Kosinus-Ähnlichkeitsmaß zurück (\textit{Top-k-Retrieval}).

\begin{equation}
  \text{sim}(\vec{q}, \vec{d}) = \frac{\vec{q} \cdot \vec{d}}{\|\vec{q}\| \cdot \|\vec{d}\|}
  \label{eq:cosine}
\end{equation}

Die abgerufenen Dokumente $D = \{d_1, \ldots, d_k\}$ werden anschließend als
zusätzlicher Kontext in den LLM-Prompt eingefügt, womit das Modell auf aktuelle,
unternehmensinterne Informationen zugreifen kann, ohne neu trainiert werden zu müssen.

\subsection{Bedeutung für das Artefakt}

Der \textit{VectorStoreNode} realisiert sowohl die Indexierungs- als auch die
Retrieval-Phase innerhalb einer einzigen n8n-Node. ChromaDB dient als persistenter
Vektorspeicher; die native SDK-Integration vermeidet den HTTP-Overhead, der bei der
Baseline-Implementierung durch drei separate n8n-Nodes entsteht (Embedding-Request,
Transformationsschritt, ChromaDB-Request).

\section{n8n im Vergleich zu anderen Low-Code-Plattformen}
\label{sec:lowcode-comparison}

\subsection{Überblick des Low-Code-Marktes}

Waszkowski (2019) definiert Low-Code-Plattformen als Systeme, die
Software\-entwicklung durch visuelle Abstraktionen für ein breites Nutzerfeld
zugänglich machen, ohne die Flexibilität für komplexe Anwendungsfälle einzuschränken
\cite{waszkowski2019lowcode}. Luo et al.\ (2021) ergänzen, dass Low-Code-Plattformen
insbesondere dann Mehrwert schaffen, wenn sie die Lücke zwischen
Business-Prozess-Experten und Softwareentwicklern schließen
\cite{luo2021characteristics}.

Im Kontext dieser Arbeit ist entscheidend, welche Plattformen den Custom-Node-Mechanismus
bereitstellen, der für die Implementierung des Artefakts erforderlich ist.
Tabelle~\ref{tab:lowcode-vergleich} vergleicht die relevantesten Plattformen.

\begin{table}[H]
  \centering
  \caption{Vergleich von Low-Code-/Workflow-Plattformen}
  \label{tab:lowcode-vergleich}
  \begin{tabularx}{\textwidth}{lXXcc}
    \toprule
    \textbf{Plattform} & \textbf{Erweiter\-barkeit} & \textbf{KI-Agent\-Support} &
    \textbf{Open Source} & \textbf{Self-hosted} \\
    \midrule
    \textbf{n8n}              & Custom Nodes (TypeScript/JS), vollständig & Nativ + erweiterbar & Ja & Ja \\
    Zapier                    & Keine eigenen Nodes; nur konfigurierbar   & Begrenzt (kein State) & Nein & Nein \\
    Make (Integromat)         & Module via Partner-\ac{API}; kein Code    & Kein nativer Support  & Nein & Nein \\
    Microsoft Power Automate  & Connectors (Low-Overhead, Enterprise)     & Copilot-Integration   & Nein & Nein \\
    Apache Airflow            & Python-Operators, vollständig             & Kein visueller Editor & Ja & Ja \\
    Node-RED                  & JS-Nodes (IoT-fokussiert)                 & Kein nativer Support  & Ja & Ja \\
    Temporal                  & Go/Java/TS-Workflows (Dev-fokussiert)     & Kein visueller Editor & Ja & Ja \\
    \bottomrule
  \end{tabularx}
\end{table}

\subsection{Detaillierter Plattformvergleich}

\subsubsection{Zapier}

Zapier ist die marktführende No-Code-Automatisierungsplattform mit über 6.000
App-Integrationen. Ihre Stärke liegt in der einfachen Verknüpfung von SaaS-Diensten
(Trigger $\to$ Action), nicht jedoch in der Implementierung benutzerdefinierter Logik.
Zapier erlaubt keine eigenen Node-Typen; Nutzer sind auf den vorgegebenen Integrations-Katalog
beschränkt. Komplexe Zustandsverwaltung ist nicht möglich \cite{luo2021characteristics}.
Für das Artefakt dieser Arbeit scheidet Zapier aus, weil die Implementierung des
ReAct-Loops benutzerdefinierte Kontrollfluss-Logik und State-Management erfordert,
die über Zapier nicht realisierbar sind.

\subsubsection{Make (ehem. Integromat)}

Make bietet einen visuellen Szenario-Editor mit Unterstützung für komplexe
Verzweigungen und Iteratoren. Eigene Module können jedoch nur über die Partner-\ac{API}
hinzugefügt werden, was erheblichen bürokratischen Aufwand bedeutet. Codebasierte
Erweiterungen sind nicht vorgesehen. Darüber hinaus ist Make ein vollständig cloudbasierter,
proprietärer Dienst ohne Self-Hosted-Option -- was für Unternehmensanwendungen mit
Datenschutz-Anforderungen problematisch sein kann \cite{luo2021characteristics}.

\subsubsection{Microsoft Power Automate}

Power Automate ist tief in das Microsoft-365-Ökosystem integriert und bietet
umfangreiche Konnektoren für Enterprise-Systeme. Die Erweiterbarkeit über Custom
Connectors ist möglich, aber auf die Microsoft-Plattform beschränkt. Eine
vollständige Self-Hosted-Deployment-Option existiert nicht. Für Unternehmen außerhalb
des Microsoft-Ökosystems entstehen erhebliche Abhängigkeiten \cite{luo2021characteristics}.

\subsubsection{Apache Airflow}

Apache Airflow ist eine auf Python basierende Workflow-Orchestrierungsplattform, die
in der Datentechnik (Data Engineering) weit verbreitet ist. Airflow ist hochgradig
erweiterbar (Custom Operators in Python), besitzt jedoch keinen visuellen
Drag-and-Drop-Editor für die Workflow-Definition. Die Zielgruppe sind Entwickler und
Data Engineers, nicht Business-Prozess-Experten. Damit erfüllt Airflow zwar die
technischen Anforderungen (State via XCom, Python-Flexibilität), nicht aber das
Transparenzversprechen eines visuellen Canvas, das Nicht-Entwicklern zugänglich ist.

\subsubsection{Node-RED}

Node-RED ist eine quelloffene, flow-basierte Entwicklungsplattform, die ursprünglich
für IoT-Anwendungen (IBM, 2013) konzipiert wurde. Die Erweiterbarkeit über JavaScript-Nodes
ist vorhanden; die Plattform ist jedoch nicht auf Enterprise-Prozessautomatisierung
ausgerichtet und bietet keinen nativen \ac{LLM}-Integration-Support. Das Ökosystem ist
deutlich kleiner als das von n8n.

\subsubsection{Temporal}

Temporal ist eine Workflow-Orchestrierungsplattform für langlebige, zustandsbehaftete
Prozesse in Go, Java oder TypeScript. Technisch wäre Temporal für die Umsetzung der
Agenten-Architektur geeignet (natives State-Management, Retry-Logic). Jedoch ist
Temporal ausschließlich entwicklerorientiert: Es gibt keinen visuellen Editor, und die
Lernkurve ist erheblich. Das Transparenz- und Zugänglichkeitsziel dieser Arbeit
wäre damit verfehlt.

\subsubsection{Warum n8n?}

n8n vereint als einzige der betrachteten Plattformen alle vier für das Artefakt
erforderlichen Eigenschaften:

\begin{enumerate}
  \item \textbf{Vollständige Erweiterbarkeit:} Custom Nodes in TypeScript/JavaScript
    ohne bürokratische Hürden, direkt per npm-Paket einbindbar.
  \item \textbf{Visueller Canvas:} Jeder Node ist ein sichtbares Element im Workflow --
    dies ist das primäre Medium für die Transparenz-Anforderung.
  \item \textbf{Self-Hosted und Open Source:} Keine Cloud-Abhängigkeit; DSGVO-konformer
    Betrieb möglich.
  \item \textbf{Natives KI-Ökosystem:} n8n bietet bereits LangChain-basierte
    KI-Nodes (\textit{AI Agent Node}), auf die das Artefakt aufbauen und die es
    übertreffen kann.
\end{enumerate}

Die Entscheidung für n8n war letztlich pragmatisch: Wer einen KI-Agenten entwickeln
will, der auch für Mitarbeitende ohne Programmierkenntnisse bedienbar und
nachvollziehbar ist, kommt an einem visuellen Workflow-Tool kaum vorbei. LangChain
ist für Entwickler gemacht -- das ist auch sein Vorteil. Aber das Ziel dieser Arbeit
war ein anderes: einen Agenten zu zeigen, dessen Entscheidungsschritte sich im Browser
verfolgen und bei Bedarf unterbrechen lassen. Das erfordert einen visuellen Canvas, und
n8n war die einzige Plattform, die das zusammen mit echter Custom-Node-Erweiterbarkeit
in TypeScript bot.

\section{Vergleich mit bestehenden KI-Agenten-Frameworks}
\label{sec:agent-comparison}

\subsection{Überblick}

Seit der Veröffentlichung von GPT-4 (OpenAI, 2023) ist eine Vielzahl von
Agenten-Frameworks entstanden, die \ac{LLM}s mit Werkzeugen verbinden. Wang et al.\
(2023) geben einen umfassenden Überblick über diese Landschaft \cite{wang2023survey}.
Für die vorliegende Arbeit sind insbesondere jene Frameworks relevant, die als
Alternativlösungen für die adressierte Problemstellung in Frage kämen.

\begin{table}[H]
  \centering
  \caption{Vergleich KI-Agenten-Frameworks mit dem Artefakt}
  \label{tab:agent-vergleich}
  \begin{tabularx}{\textwidth}{lXcccc}
    \toprule
    \textbf{Framework} & \textbf{Paradigma} & \textbf{Visuell} & \textbf{Transparent} &
    \textbf{HITL} & \textbf{No-Code} \\
    \midrule
    \textbf{Artefakt (n8n)} & ReAct, 1 Step/Node & \checkmark & \checkmark & \checkmark & \checkmark \\
    n8n AI Agent Node       & LangChain-intern   & \checkmark & \texttimes & Begrenzt & \checkmark \\
    LangChain (Python)      & Chain/Agent        & \texttimes & \texttimes & Manuell  & \texttimes \\
    AutoGPT                 & Autonome Schleife  & \texttimes & \texttimes & \texttimes & \texttimes \\
    Claude Code             & Coding-Agent       & \texttimes & \texttimes & Teilweise & \texttimes \\
    CrewAI                  & Multi-Agent        & \texttimes & Teilweise  & \texttimes & \texttimes \\
    AutoGen (Microsoft)     & Multi-Agent, Chat  & \texttimes & Teilweise  & Teilweise & \texttimes \\
    LlamaIndex              & RAG-fokussiert     & \texttimes & \texttimes & \texttimes & \texttimes \\
    \bottomrule
  \end{tabularx}
\end{table}

\subsection{n8n AI Agent Node (Baseline)}

Der native \textit{AI Agent Node} von n8n basiert intern auf dem LangChain-Framework
und bietet eine visuelle Konfiguration. Er hat jedoch drei strukturelle Einschränkungen,
die das Artefakt behebt:

\begin{itemize}
  \item \textbf{Keine persistente Zustandsverwaltung:} Der AI Agent Node führt die
    gesamte ReAct-Schleife intern aus. Der Zwischenzustand ist nicht im Canvas sichtbar
    und kann nicht zwischen Node-Ausführungen persistiert werden.
  \item \textbf{Eingeschränkte Transparenz:} Die Thought-Action-Observation-Tripel
    sind zwar in den Execution-Logs abrufbar, aber nicht als diskrete, interagierbare
    Workflow-Elemente visualisiert.
  \item \textbf{Kein granulares Human-in-the-Loop:} Der Node pausiert nicht nach jeder
    Action; menschliche Eingriffe sind nur am Anfang oder Ende des Prozesses möglich.
\end{itemize}

\subsection{LangChain}

LangChain ist das meistgenutzte Python-Framework für \ac{LLM}-Anwendungen. Es bietet
umfangreiche Abstraktionen für Chains, Agenten, Retriever und Tools. Für die
Implementierung eines zustandsbehafteten ReAct-Agenten ist LangChain technisch
geeignet. Die Einschränkungen im Kontext dieser Arbeit:

\begin{itemize}
  \item \textbf{Kein visueller Canvas:} Die gesamte Logik ist in Python-Code
    definiert. Business-Prozess-Experten ohne Python-Kenntnisse können den Agenten
    weder konfigurieren noch seinen Ablauf nachvollziehen.
  \item \textbf{Erschwertes Human-in-the-Loop:} LangChain bietet Callbacks und
    \textit{LangGraph} für HITL-Szenarien, jedoch ohne visuelle Repräsentation.
  \item \textbf{Schlechte Integration in bestehende IT-Infrastruktur:} Die Integration
    in ERP-Systeme, SQL-Datenbanken und REST-\ac{API}s erfordert manuelle Python-Entwicklung.
    n8n's umfangreiche Node-Bibliothek (400+ Integrationen) steht nicht zur Verfügung.
\end{itemize}

\subsection{AutoGPT}

AutoGPT war eines der ersten öffentlich verfügbaren autonomen Agenten-Systeme. Es
iteriert selbstständig über ReAct-ähnliche Zyklen bis zur Aufgabenerfüllung und benötigt
dabei keine menschlichen Eingriffe. Genau darin liegt das zentrale Problem für
produktive Unternehmensanwendungen: AutoGPT kann unkontrolliert irreversible Aktionen
ausführen -- etwa Dateien löschen, E-Mails versenden oder Web-\ac{API}s aufrufen --,
ohne dass ein Operator eingreifen kann. Transparenz ist nicht vorgesehen; der
Entscheidungsprozess ist vollständig in einer Black Box verborgen.

\subsection{Claude Code}

Claude Code (Anthropic, 2024) ist ein KI-gestützter Coding-Assistent, der als
Command-Line-Interface direkt im Terminal läuft und Entwicklern beim Schreiben,
Refactoring und Debuggen von Code hilft. Für den Kontext der vorliegenden Arbeit
ist Claude Code aus zwei Gründen relevant, aber nicht als Lösung geeignet:

\begin{itemize}
  \item \textbf{Domänenspezifisch:} Claude Code ist auf Software-Entwicklungsaufgaben
    spezialisiert (Datei-Operationen, Git, Terminal). Die in dieser Arbeit adressierten
    Geschäftsprozess-Szenarien (ERP, SQL, RAG) sind nicht der Anwendungsfall.
  \item \textbf{Keine Prozessintegration:} Claude Code besitzt keinen visuellen
    Workflow-Editor und keine nativen Konnektoren zu Unternehmensanwendungen. Jede
    Integration müsste manuell programmiert werden.
  \item \textbf{Black-Box-Entscheidungen:} Ähnlich wie AutoGPT sind die internen
    Reasoning-Schritte nicht für externe Auditoren sichtbar.
\end{itemize}

Claude Code richtet sich in erster Linie an Entwickler, die ihre Werkzeuge über
Code steuern wollen und mit der zugrunde liegenden Technologie vertraut sind.
Das in dieser Arbeit entwickelte Artefakt verfolgt einen anderen Ansatz: Es soll
auch für Personen ohne Programmierkenntnisse nutzbar sein und den Ablauf des
Agenten zu jedem Zeitpunkt transparent und nachvollziehbar halten.

\subsection{CrewAI und AutoGen}

CrewAI und AutoGen (Microsoft) realisieren Multi-Agenten-Architekturen, in denen
mehrere spezialisierte Agenten miteinander kommunizieren \cite{wang2023survey}. Beide
Frameworks sind technisch anspruchsvoll und Python-basiert. Für die vorliegende Arbeit
-- die einen einzelnen ReAct-Agenten in einer Enterprise-Low-Code-Umgebung verankert
-- sind Multi-Agenten-Frameworks überengineered und erhöhen die Komplexität ohne
direkten Mehrwert.

\subsection{LlamaIndex}

LlamaIndex (früher GPT Index) ist ein \ac{RAG}-fokussiertes Python-Framework, das
umfangreiche Dokumenten-Indexierungspipelines und Retrieval-Strategien bietet.
Technisch ist LlamaIndex für den VectorStore-Anwendungsfall gut geeignet. Es fehlt
jedoch der visuelle Workflow-Editor und die Möglichkeit, den Retrieval-Prozess ohne
Programmierkenntnisse zu konfigurieren.

\section{Design Science Research als Forschungsmethodik}
\label{sec:dsr}

Hevner et al.\ (2004) definieren \ac{DSR} als wissenschaftliche Methode, bei der
IT-Artefakte -- Konstrukte, Modelle, Methoden oder Instanziierungen -- als
Primärergebnis der Forschung entwickelt und evaluiert werden \cite{hevner2004design}.
Die Methodik kombiniert die rigorose Evaluation aus der Naturwissenschaft mit dem
konstruktiven Ansatz aus der Ingenieurwissenschaft.

Für die vorliegende Arbeit bedeutet dies:

\begin{itemize}
  \item \textbf{Artefakt:} Das n8n-Agentic-Pack (3 Custom Nodes, Evaluation-Framework,
    Dokumentation).
  \item \textbf{Problemraum:} Fehlende Transparenz und Zustandslosigkeit von LLM-Agenten
    in Low-Code-Plattformen.
  \item \textbf{Wissensraum:} Einschlägige Publikationen zu ReAct, RAG, Memory, DSR.
  \item \textbf{Evaluation:} Drei Benchmark-Szenarien mit quantitativen und qualitativen
    Metriken, verglichen mit Baseline und Referenz.
\end{itemize}

% =============================================================================
% KAPITEL 3: AGENTSTATETNODE
% =============================================================================
\chapter{Implementierung: AgentStateNode}

\section{Theoretische Grundlage}

Der \textit{AgentStateNode} implementiert das persistente Zustandsmodell für den
gesamten Agenten-Lebenszyklus. Theoretische Basis ist der Memory-Stream-Ansatz aus
Park et al.\ (2023), ergänzt durch die Sicherheitsmechanismen aus Wang et al.\ (2023)
\cite{park2023generative, wang2023survey}.

Die zentrale Designentscheidung besteht darin, den Agentenzustand nicht als internes
Objekt innerhalb eines Agenten-Frameworks (wie LangChain's \texttt{AgentExecutor}) zu
verwalten, sondern ihn als explizites \ac{JSON}-Objekt im n8n-Datenstrom zu
materialisieren. Dieser Ansatz ermöglicht:

\begin{itemize}
  \item Sichtbarkeit des Zustands nach jeder Node-Ausführung im n8n-Canvas
  \item Persistenz des Zustands in externen Speichern (Redis, PostgreSQL) ohne
    Änderung der Node-Implementierung
  \item Reproduzierbarkeit: Ein gespeicherter Zustand kann jederzeit geladen und
    fortgesetzt werden
\end{itemize}

\section{Designentscheidungen und Alternativen}

\subsection{Entscheidung 1: Strukturierter History-Log vs. Freitext-Kontext}

\paragraph{Gewählter Ansatz:}
Der Agentenzustand speichert die vergangenen Schritte als typisiertes Array von
\texttt{AgentStep}-Objekten (je Thought, Action, ActionInput, Observation, Timestamp).

\paragraph{Alternative A -- Freitext-Kontext:}
Der gesamte Gesprächsverlauf wird als einfacher String gespeichert (ähnlich wie
LangChain's \texttt{ConversationBufferMemory}). Vorteil: geringerer Implementierungsaufwand.
Nachteil: kein maschineller Zugriff auf einzelne Felder; Benchmarking-Auswertung
(Metrik: \texttt{steps\_to\_completion}) erfordert aufwändiges Text-Parsing.

\paragraph{Alternative B -- Ring Buffer (FIFO):}
Nur die letzten $n$ Schritte werden gespeichert (ähnlich
\texttt{ConversationBufferWindowMemory} in LangChain). Vorteil: begrenzter
Speicherbedarf, auch bei langen Agenten-Läufen. Nachteil: Schritte gehen verloren;
der Agent kann frühere Beobachtungen nicht mehr referenzieren. Dies ist insbesondere
für den SC-03-Benchmark (Error Recovery) kritisch, wo der Agent auf den fehlgeschlagenen
ersten Schritt zurückblicken muss.

\paragraph{Begründung der Wahl:}
Der vollständige, strukturierte History-Log wurde gewählt, weil (1) die wissenschaftliche
Auswertung Zugriff auf alle Schritte benötigt und (2) Park et al.\ (2023) explizit
zeigen, dass der vollständige Memory Stream -- nicht ein abgeschnittener Kontext --
zu kohärentem Langzeitverhalten führt \cite{park2023generative}.

\subsection{Entscheidung 2: Serialisierung als JSON im n8n-Datenstrom}

\paragraph{Gewählter Ansatz:}
Der Zustand wird als \ac{JSON}-Objekt zwischen Nodes weitergereicht, was dem nativen
Datenmodell von n8n entspricht.

\paragraph{Alternative A -- Redis / externer KV-Speicher:}
Der Zustand wird unter einer \texttt{sessionId} in Redis gespeichert; Nodes lesen
und schreiben ihn über eine \ac{API}. Vorteil: Zustand ist von der Workflow-Instanz
entkoppelt. Nachteil: zusätzliche Infrastrukturabhängigkeit; nicht ohne Weiteres
in Docker-Compose-Testumgebungen integrierbar.

\paragraph{Alternative B -- n8n Static Data:}
n8n bietet einen \texttt{staticData}-Mechanismus für persistente Node-Zustände.
Nachteil: nicht für Multi-Instance-Deployments geeignet; nicht skalierbar.

\paragraph{Begründung der Wahl:}
Die JSON-Serialisierung im Datenstrom wurde gewählt, weil sie ohne zusätzliche
Infrastruktur auskommt und das Artefakt so in jeder n8n-Umgebung ohne Vorkonfiguration
lauffähig ist. Für produktive Deployments kann der Zustand trivial in einen externen
Speicher ausgelagert werden, ohne die Node-Implementierung zu ändern.

\subsection{Entscheidung 3: Harte Iterationsgrenze (maxSteps)}

\paragraph{Gewählter Ansatz:}
Der Parameter \texttt{maxSteps} wirkt als harte Grenze; das Überschreiten löst einen
\texttt{NodeOperationError} aus.

\paragraph{Alternative -- Weiche Warnung:}
Eine Warnung wird in das Observations-Feld geschrieben; der Agent kann entscheiden,
ob er trotzdem fortfährt. Nachteil: ein fehlerhafter Agent ignoriert möglicherweise
die Warnung und erzeugt unkontrollierte \ac{API}-Kosten.

\paragraph{Begründung:}
Wang et al.\ (2023) identifizieren Endlosschleifen als kritischen Ausfallmodus in
Produktivsystemen \cite{wang2023survey}. Eine harte Grenze ist die robustere Wahl.

\section{Zustandsschema}

\begin{lstlisting}[language=Java, caption={AgentState-Schnittstelle in TypeScript},
  label=lst:agentstate-schema]
interface AgentState {
  sessionId: string;          // UUID-v4: eindeutiger Session-Schluessel
  context: string;            // Aufgabenbeschreibung im natuerlichen Sprachformat
  history: AgentStep[];       // Memory Stream (append-only, vollstaendig)
  stepCount: number;          // Anzahl abgeschlossener Schritte
  maxSteps: number;           // Harte Iterationsgrenze
  metadata: Record<string, unknown>; // Erweiterungsdaten
}

interface AgentStep {
  step: number;       // 1-basierter Sequenzindex
  thought: string;    // Reasoning-Spur des LLM
  action: string;     // Ausgewaehltes Tool
  actionInput: Record<string, unknown>; // JSON-Parameter fuer das Tool
  observation: string;// Tool-Ergebnis
  timestamp: string;  // ISO 8601
}
\end{lstlisting}

\section{Implementierte Operationen}

Der \textit{AgentStateNode} stellt vier Operationen bereit, die den gesamten
Lebenszyklus eines Agenten-Runs abdecken:

\begin{description}[leftmargin=3cm]
  \item[\texttt{initialize}] Erstellt einen neuen Agentenzustand mit einer
    frischen UUID-v4-Session-ID, einem leeren History-Array und dem angegebenen
    Kontext. Entspricht dem \textit{Bootstrapping}-Schritt eines generativen Agenten
    (Park et al., 2023).

  \item[\texttt{update}] Hängt einen abgeschlossenen \texttt{AgentStep} an den
    History-Array und inkrementiert \texttt{stepCount}. Vor der Ausführung wird
    geprüft, ob \texttt{maxSteps} erreicht ist.

  \item[\texttt{get}] Gibt den aktuellen Zustand nach Validierung aller Pflichtfelder
    zurück. Dient zur Inspektion des Zustands an beliebiger Stelle im Workflow.

  \item[\texttt{reset}] Löscht \texttt{history} und \texttt{stepCount}, behält jedoch
    \texttt{sessionId}, \texttt{context}, \texttt{maxSteps} und \texttt{metadata}.
    Ermöglicht einen Neustart ohne neue Session-Erstellung (relevant für SC-03).
\end{description}

\section{Reine Funktionen und Testbarkeit}

Alle sechs reinen Geschäftslogik-Funktionen (\texttt{initializeState},
\texttt{appendStep}, \texttt{isMaxStepsReached}, \texttt{resetState},
\texttt{serializeState}, \texttt{deserializeState}) sind unabhängig von der n8n-Node-Klasse
implementiert und vollständig unit-testbar. Diese Trennung folgt dem
\textit{Separation-of-Concerns}-Prinzip und erlaubt es, die Kernlogik zu testen,
ohne n8n-Infrastruktur starten zu müssen.

\section{Testergebnisse}

Die Testsuite \texttt{nodes/AgentStateNode/\_\_tests\_\_/AgentStateNode.test.ts}
umfasst 40 Testfälle in sieben Testsuiten. Die Jest-Coverage-Auswertung ergibt:

\begin{table}[H]
  \centering
  \caption{Testabdeckung AgentStateNode}
  \label{tab:coverage-agentstate}
  \begin{tabular}{lrrrr}
    \toprule
    \textbf{Metrik} & \textbf{Statements} & \textbf{Branches} & \textbf{Functions} & \textbf{Lines} \\
    \midrule
    Ergebnis & 100\,\% & 83,87\,\% & 100\,\% & 100\,\% \\
    Zielgröße & $\geq 80\,\%$ & $\geq 75\,\%$ & $\geq 80\,\%$ & $\geq 80\,\%$ \\
    \bottomrule
  \end{tabular}
\end{table}

\section{Bezug zur Forschungsfrage}

Der \textit{AgentStateNode} beantwortet \textbf{TF1} (Machbarkeit), indem er zeigt,
dass persistente Zustandsverwaltung innerhalb des n8n-Custom-Node-Mechanismus vollständig
realisierbar ist. Er adressiert \textbf{TF2} (Transparenz), weil der gesamte
Entscheidungsverlauf nach jeder Ausführung als strukturiertes JSON-Objekt im Canvas
sichtbar ist -- ein qualitativer Vorteil, den die Baseline-Implementierung nicht bietet.

% =============================================================================
% KAPITEL 4: VECTORSTORENODE
% =============================================================================
\chapter{Implementierung: VectorStoreNode}

\section{Theoretische Grundlage}

Der \textit{VectorStoreNode} realisiert die Retriever-Komponente des \ac{RAG}-Frameworks
(Lewis et al., 2020) als nativen n8n-Node \cite{lewis2020retrieval}. Die zentrale
Motivation ist die Eliminierung des in der Baseline-Implementierung erforderlichen
HTTP-Workarounds, der drei separate n8n-Nodes (Embedding-Request, Transformations-Node,
ChromaDB-Request) benötigt.

ChromaDB wurde als Vektorspeicher gewählt, weil es:
\begin{enumerate}
  \item vollständig quelloffen und selbst-hostbar ist (DSGVO-Kompatibilität)
  \item eine native TypeScript/JavaScript-Client-Bibliothek (\texttt{chromadb} npm)
    bereitstellt
  \item persistenten, containerisierten Betrieb via Docker unterstützt
  \item alle gängigen Ähnlichkeitsmaße unterstützt (Kosinus, L2, Inneres Produkt)
\end{enumerate}

\section{Designentscheidungen und Alternativen}

\subsection{Entscheidung 1: ChromaDB vs. andere Vektorspeicher}

\begin{table}[H]
  \centering
  \caption{Vergleich von Vektorspeicher-Optionen}
  \label{tab:vectordb-vergleich}
  \begin{tabularx}{\textwidth}{lXccc}
    \toprule
    \textbf{Speicher} & \textbf{Stärken} & \textbf{Self-hosted} & \textbf{TS-SDK} & \textbf{Open Source} \\
    \midrule
    \textbf{ChromaDB}  & Einfach, Docker-ready, akademisch weit verbreitet & Ja & Ja & Ja \\
    Pinecone           & Hochskalierbar, managed                           & Nein & Ja & Nein \\
    Weaviate           & GraphQL, umfangreich                              & Ja & Ja & Ja \\
    pgvector           & PostgreSQL-Extension, keine neue Infrastruktur     & Ja & Ja & Ja \\
    Qdrant             & Performant, Rust-basiert                          & Ja & Ja & Ja \\
    \bottomrule
  \end{tabularx}
\end{table}

\paragraph{Pinecone} ist ein managed Vektordatenbankdienst ohne Self-Hosted-Option.
Für eine Bachelorarbeit, die eine reproduzierbare Docker-Compose-Testumgebung
bereitstellen soll, scheidet Pinecone aus.

\paragraph{pgvector} ist eine PostgreSQL-Extension und hätte den Vorteil, keine
neue Datenbankinfrastruktur zu erfordern. Nachteil: eine PostgreSQL-Instanz muss
bereits existieren; der Konfigurationsaufwand ist höher als bei ChromaDB.

\paragraph{ChromaDB} wurde gewählt, weil es den geringsten Setup-Aufwand bietet
(ein Docker-Container), über eine vollständige TypeScript-Bibliothek verfügt und
in der akademischen Community für \ac{RAG}-Prototypen weit verbreitet ist.

\subsection{Entscheidung 2: Native SDK vs. HTTP-Request}

\paragraph{Gewählter Ansatz:}
Direkte Nutzung der \texttt{chromadb}-npm-Bibliothek innerhalb des Custom Node.

\paragraph{Alternative -- HTTP-Request-Node (Baseline):}
Die Standard-n8n-Implementierung nutzt den generischen HTTP-Request-Node für alle
ChromaDB-Aufrufe. Dies erfordert für eine Upsert-Operation:
\begin{enumerate}
  \item HTTP POST \texttt{/v1/embeddings} (OpenAI-\ac{API}) $\to$ Embedding-Vektor
  \item Set-Node: Manuelles Mapping des Vektors in das ChromaDB-Format
  \item HTTP POST \texttt{/collections/\{name\}/upsert} (ChromaDB-\ac{API})
\end{enumerate}

\paragraph{Gemessene Vorteile der SDK-Integration:}
\begin{itemize}
  \item \textbf{Node-Reduktion:} 3 Nodes $\to$ 1 Node pro Operation
  \item \textbf{Fehlervermeidung:} Manuelle JSON-Transformation in Schritt 2 ist
    fehleranfällig (falsche Feldnamen, fehlende Felder)
  \item \textbf{Latenz:} Kein HTTP-Parse-Overhead zwischen den drei Aufrufen; die
    SDK-Aufrufe laufen im selben Node-Prozess
  \item \textbf{Typsicherheit:} TypeScript-Typen der SDK verhindern Laufzeitfehler
\end{itemize}

\subsection{Entscheidung 3: Upsert-Semantik}

Dokumente werden per \texttt{upsert} (create-or-update) statt per \texttt{insert}
gespeichert. Dies entspricht der Empfehlung von Lewis et al.\ (2020) für produktive
\ac{RAG}-Systeme, bei denen Dokumente regelmäßig aktualisiert werden \cite{lewis2020retrieval}.
Ein reines Insert würde bei wiederholten Indexierungsläufen Duplikate erzeugen.

\subsection{Entscheidung 4: Austauschbares Embedding-Modell}

Das Embedding-Modell ist als Dropdown-Parameter konfigurierbar. Drei OpenAI-Modelle
werden nativ unterstützt (\texttt{text-embedding-3-small}, \texttt{text-embedding-3-large}, \\
\texttt{text-embedding-ada-002}). Die \texttt{embedText}-Funktion ist bewusst
isoliert, um zukünftige lokale Modelle (z.\,B. Ollama/nomic-embed) ohne Änderung der
Speicherlogik einzubinden.

\section{Implementierte Operationen}

\begin{description}[leftmargin=3cm]
  \item[\texttt{upsert}] Generiert das Embedding für den übergebenen Text und
    speichert Vektor, Rohtext und Metadaten in der angegebenen ChromaDB-Kollektion.

  \item[\texttt{query}] Generiert das Embedding für die Suchanfrage und gibt die
    $k$ ähnlichsten Dokumente mit \texttt{id}, \texttt{document}, \texttt{metadata}
    und \texttt{distance} zurück.

  \item[\texttt{delete}] Entfernt ein Dokument anhand seiner ID aus der Kollektion.

  \item[\texttt{listCollections}] Listet alle vorhandenen Kollektionen; Hilfsfunktion
    für Debugging und Workflow-Entwicklung.
\end{description}

\section{Testergebnisse}

\begin{table}[H]
  \centering
  \caption{Testabdeckung VectorStoreNode}
  \label{tab:coverage-vector}
  \begin{tabular}{lrrrr}
    \toprule
    \textbf{Metrik} & \textbf{Statements} & \textbf{Branches} & \textbf{Functions} & \textbf{Lines} \\
    \midrule
    Ergebnis & 100\,\% & 77,77\,\% & 100\,\% & 100\,\% \\
    Zielgröße & $\geq 80\,\%$ & $\geq 75\,\%$ & $\geq 80\,\%$ & $\geq 80\,\%$ \\
    \bottomrule
  \end{tabular}
\end{table}

Alle 29 Testfälle verwenden Jest-Mocks für OpenAI-SDK und ChromaDB-Client, um
offline-fähige, deterministische Tests zu gewährleisten.

\section{Bezug zur Forschungsfrage}

Der \textit{VectorStoreNode} adressiert \textbf{TF3} (Integration) direkt: Die
native SDK-Integration reduziert die Node-Anzahl für eine Upsert-/Query-Operation
von drei auf eins. Die Benchmark-Szenarien SC-01 und SC-02 messen diese Reduktion
quantitativ über die Metrik \texttt{node\_count}.

% =============================================================================
% KAPITEL 5: REACTLOOPNODE
% =============================================================================
\chapter{Implementierung: ReActLoopNode}

\section{Theoretische Grundlage}

Der \textit{ReActLoopNode} ist die orchestrierende Kernkomponente des Artefakts.
Er implementiert den \ac{TAO}-Zyklus des ReAct-Musters (Yao et al., 2022) als
diskreten, schrittweisen n8n-Node \cite{yao2022react}. Die Werkzeugbeschreibungen
im Prompt folgen dem Toolformer-Paradigma (Schick et al., 2023), bei dem Tools durch
Namen und natürlichsprachliche Beschreibungen für das \ac{LLM} zugänglich gemacht
werden, ohne explizites Fine-Tuning zu erfordern \cite{schick2023toolformer}.

\section{Designentscheidungen und Alternativen}

\subsection{Entscheidung 1: Ein Schritt pro Node-Ausführung (Single-Step-Design)}

\paragraph{Gewählter Ansatz:}
Der \textit{ReActLoopNode} führt \emph{genau einen} ReAct-Schritt pro Node-Ausführung
durch: Er baut den Prompt auf, ruft das \ac{LLM} auf, parst die Ausgabe und gibt das
Ergebnis an den n8n-Workflow zurück. Die eigentliche Schleife entsteht durch einen
Feedback-Pfeil im n8n-Canvas zurück zum Node.

\paragraph{Alternative -- Internal-Loop-Design:}
Der Node führt alle Schritte bis zur Final Answer intern in einer \texttt{while}-Schleife
aus. Vorteile: einfachere Workflow-Struktur; weniger Canvas-Objekte. Nachteile:
\begin{itemize}
  \item Der gesamte Entscheidungsprozess läuft unsichtbar in der Black Box des Nodes.
  \item Kein Human-in-the-Loop zwischen Schritten möglich.
  \item n8n's native Execution-Monitoring kann keine Zwischen\-stände anzeigen.
  \item Entspricht dem Muster des Baseline-Nodes (AI Agent Node) -- der gerade
    überwunden werden soll.
\end{itemize}

\paragraph{Begründung:}
Das Single-Step-Design ist die direkte Umsetzung des Transparenzziels aus der
Forschungsfrage. Amershi et al.\ (2019) formulieren in Guideline H7 (Human-AI
Interaction): \textit{``Support efficient invocation -- make it easy for people
to invoke or dismiss the AI system.''} \cite{amershi2019guidelines}. Das duale
Ausgabe-Routing (Action Required / Final Answer) realisiert genau dieses Prinzip:
Jeder Schritt ist ein Checkpoint, an dem ein Operator entscheiden kann.

\subsection{Entscheidung 2: Duales Ausgabe-Routing}

\paragraph{Gewählter Ansatz:}
Der Node besitzt zwei Ausgänge:
\begin{itemize}
  \item \textbf{Output 0 -- Action Required:} Aktiviert, wenn das \ac{LLM} ein
    weiteres Tool aufrufen möchte.
  \item \textbf{Output 1 -- Final Answer:} Aktiviert, wenn das \ac{LLM} die Aufgabe
    als abgeschlossen erklärt.
\end{itemize}

\paragraph{Alternative -- Single-Output mit Typ-Flag:}
Ein einzelner Ausgangspfad mit einem \texttt{isFinalAnswer}-Boolefeld; ein
nachgelagerter If-Node routet. Nachteil: die Routing-Logik liegt außerhalb des
ReAct-Nodes und muss vom Workflow-Entwickler manuell konfiguriert werden.

\paragraph{Begründung:}
Das duale Routing kapsel die Terminierungslogik im Node und vereinfacht den
Workflow-Canvas erheblich.

\subsection{Entscheidung 3: Text-basiertes Prompt-Format vs. Function Calling}

\paragraph{Gewählter Ansatz:}
Das ReAct-Prompt-Template schreibt ein striktes Zeilenformat vor (\texttt{Thought:},
\texttt{Action:}, \texttt{Action Input:}, \texttt{Final Answer:}), das per
Zeilenpräfix-Matching geparst wird.

\paragraph{Alternative -- OpenAI Function Calling:}
Neuere \ac{LLM}s (GPT-4o, Claude 3) unterstützen natives Function Calling via
\ac{API}, bei dem das Modell strukturierte JSON-Ausgaben für Werkzeugaufrufe
produziert. Vorteile: zuverlässigeres Parsing, keine Prompt-Formatierungsregeln
notwendig. Nachteile:
\begin{itemize}
  \item Modellspezifisch: Nicht alle \ac{LLM}s unterstützen Function Calling.
  \item Weniger transparent: Der \textit{Thought} (Reasoning-Spur) ist beim
    Function Calling nicht Bestandteil der Standardausgabe.
\end{itemize}

\paragraph{Begründung:}
Das textbasierte Format wurde gewählt, weil es modellunabhängig ist und den
\textit{Thought}-Schritt explizit erzwingt. Der Thought ist das zentrale
Transparenz-Element des ReAct-Musters; ohne ihn wäre der Agenten-Prozess nicht
besser auditierbar als die Baseline \cite{yao2022react}.

\subsection{Entscheidung 4: Temperature = 0 als Standard}

Für Tool-Use-Szenarien empfehlen Yao et al.\ (2022) niedrige Temperaturen, da
deterministische Ausgaben das Parsing der \texttt{Action}/\texttt{Action Input}-Felder
zuverlässiger machen \cite{yao2022react}. Der Standard-Wert von 0,0 kann vom Nutzer
erhöht werden, falls kreative Aufgaben bearbeitet werden sollen.

\section{Prompt-Aufbau}

Das Prompt-Template wird durch \texttt{buildReActPrompt(state, tools)} erzeugt.
Es enthält vier strukturierte Abschnitte:

\begin{enumerate}
  \item \textbf{Context:} Die natürlichsprachliche Aufgabenbeschreibung aus dem
    \textit{AgentState}.
  \item \textbf{Available Tools:} Liste aller Tools als \textit{Name: Beschreibung}-Paare
    (Toolformer-Format, Schick et al., 2023).
  \item \textbf{Previous Steps:} Alle bisherigen \ac{TAO}-Tripel aus dem
    \textit{Memory Stream} (Park et al., 2023).
  \item \textbf{Output-Anweisungen:} Striktes Format für Thought/Action/Final Answer.
\end{enumerate}

\section{Antwort-Parsing}

Die Funktion \texttt{parseReActResponse(raw)} implementiert einen einfachen
Zeilenpräfix-Parser ohne externe Bibliotheken. Abgedeckte Szenarien:
\begin{itemize}
  \item Normale Action-Antwort (Thought + Action + Action Input)
  \item Final-Answer-Antwort (Thought + Final Answer)
  \item Nicht-JSON-Action-Input (wird als \texttt{\{raw: "..."\}} gekapselt)
  \item Fehlende Felder (Fehlerbehandlung mit sprechender Fehlermeldung)
\end{itemize}

\section{Human-in-the-Loop}

Das \texttt{humanInTheLoop}-Flag (Boolean, Standard: \texttt{false}) versieht die
Ausgabe mit einem \texttt{awaitingApproval}-Kennzeichen. In Verbindung mit einem
n8n-\textit{Wait}-Node kann der Workflow so konfiguriert werden, dass er nach jeder
Action auf eine manuelle Freigabe wartet, bevor der Tool-Node ausgeführt wird.
Amershi et al.\ (2019) betonen, dass HITL-Mechanismen ein zentrales Qualitätsmerkmal
menschenzentrierter KI-Systeme sind \cite{amershi2019guidelines}.

\section{Testergebnisse}

\begin{table}[H]
  \centering
  \caption{Testabdeckung ReActLoopNode}
  \label{tab:coverage-react}
  \begin{tabular}{lrrrr}
    \toprule
    \textbf{Metrik} & \textbf{Statements} & \textbf{Branches} & \textbf{Functions} & \textbf{Lines} \\
    \midrule
    Ergebnis & 100\,\% & 82,05\,\% & 100\,\% & 100\,\% \\
    Zielgröße & $\geq 80\,\%$ & $\geq 75\,\%$ & $\geq 80\,\%$ & $\geq 80\,\%$ \\
    \bottomrule
  \end{tabular}
\end{table}

\section{Bezug zur Forschungsfrage}

Der \textit{ReActLoopNode} ist die direkte Antwort auf \textbf{TF2} (Transparenz):
Jeder \ac{TAO}-Zyklus wird als diskreter Node-Durchlauf im Canvas sichtbar. Das
HITL-Flag adressiert die Kontrollierbarkeits-Komponente der Forschungsfrage; die
\texttt{maxIterations}-Grenze verhindert unkontrollierte Ausführungen.

% =============================================================================
% KAPITEL 6: EVALUATION
% =============================================================================
\chapter{Evaluation}

\section{Evaluationsrahmen nach Design Science Research}

Die Evaluation folgt dem DSR-Evaluationsrahmen nach Hevner et al.\ (2004), der den
Vergleich des Artefakts mit einer Baseline und einer Referenz-Implementierung einfordert
\cite{hevner2004design}. Drei Implementierungen werden untersucht:

\begin{table}[H]
  \centering
  \caption{Evaluierte Implementierungen}
  \label{tab:eval-implementations}
  \begin{tabularx}{\textwidth}{llX}
    \toprule
    \textbf{Bezeichnung} & \textbf{System} & \textbf{Beschreibung} \\
    \midrule
    Baseline  & Standard n8n  & n8n AI Agent Node + HTTP-Request-Nodes für ChromaDB \\
    Artefakt  & Agentic Pack  & Die drei Custom Nodes dieser Arbeit \\
    Referenz  & LangChain     & LangChain ReAct Agent (Python), außerhalb n8n \\
    \bottomrule
  \end{tabularx}
\end{table}

\section{Benchmark-Szenarien}

\subsection{SC-01: Data Pipeline}

\paragraph{Beschreibung:}
Der Agent ruft alle ausstehenden Bestellungen (\texttt{status='pending'}) aus einer
simulierten SQL-Datenbank ab, transformiert sie in das ERP-Format (Felder:
\texttt{order\_id}, \texttt{customer\_name}, \texttt{total\_eur}) und überträgt das
Ergebnis an das ERP-System.

\paragraph{Relevanz:}
SC-01 testet die Integrationstiefe und Tool-Verkettungsfähigkeit. Autonome
Agenten-Frameworks (LangChain, AutoGPT) verfügen nicht über native ERP- oder
Datenbankverbindungen; die Hypothese ist, dass das Artefakt hier einen messbaren
Vorteil bei \texttt{integration\_score} und \texttt{node\_count} aufweist.

\paragraph{Erwarteter Ablauf:}
\texttt{sql\_query} $\to$ \texttt{transform\_data} $\to$ \texttt{erp\_push} $\to$ Final Answer
(3 Schritte)

\subsection{SC-02: RAG Question Answering}

\paragraph{Beschreibung:}
Der Agent beantwortet eine Frage über interne Unternehmensdokumente (\textit{Welchen
Gesamtumsatz erzielte Acme GmbH in Q2 2024?}) mittels semantischer Suche im
VectorStoreNode.

\paragraph{Relevanz:}
SC-02 misst die RAG-Effizienz. Die Hypothese ist, dass der native \textit{VectorStoreNode}
niedrigere \texttt{latency\_ms} und \texttt{node\_count} aufweist als die
HTTP-Baseline (3 Nodes $\to$ 1 Node). Zusätzlich werden RAG-spezifische Metriken
erhoben (Precision@1, Precision@3).

\subsection{SC-03: Error Recovery}

\paragraph{Beschreibung:}
In Schritt 1 wird ein injizierter Datenbankfehler (Timeout) ausgelöst. Der Agent
muss den Fehler erkennen, eine Fallback-Strategie wählen und die Aufgabe (Top-5
überfällige Rechnungen ins ERP) trotzdem abschließen.

\paragraph{Relevanz:}
SC-03 testet die Zustandspersistenz über Fehler hinweg. In der Baseline-Implementierung
gibt es keine persistente Zustandsverwaltung; ein Fehler erzwingt einen kompletten
Neustart. Das Artefakt kann den Zustand nach dem Fehler laden und direkt fortsetzen.
Die Hypothese lautet: \texttt{error\_recovery\_success = true} für das Artefakt,
\texttt{false} für die Baseline.

\section{Erhobene Metriken}

\subsection{Quantitative Metriken}

\begin{table}[H]
  \centering
  \caption{Quantitative Benchmark-Metriken}
  \label{tab:metrics-quant}
  \begin{tabularx}{\textwidth}{llX}
    \toprule
    \textbf{Metrik} & \textbf{Typ} & \textbf{Beschreibung} \\
    \midrule
    \texttt{token\_count\_total}    & int  & Summe aller Input- und Output-Tokens \\
    \texttt{node\_count}            & int  & Anzahl der Nodes im n8n-Workflow \\
    \texttt{steps\_to\_completion}  & int  & Anzahl ReAct-Iterationen bis Final Answer \\
    \texttt{success}                & bool & Aufgabe korrekt abgeschlossen \\
    \texttt{error\_recovery\_success} & bool & Fehler-Recovery erfolgreich (nur SC-03) \\
    \texttt{latency\_ms}            & int  & Wanduhrzeit in Millisekunden \\
    \bottomrule
  \end{tabularx}
\end{table}

\subsection{Qualitative Metriken}
\label{sec:qualitative-metriken}

Die qualitativen Metriken werden vom Evaluator (Thesis-Autor) anhand eines
standardisierten Scoring-Bogens bewertet (Skala 1--5), der auf den
Human-AI-Interaction-Guidelines von Amershi et al.\ (2019) basiert \cite{amershi2019guidelines}.
Im Unterschied zu den quantitativen Metriken weisen die qualitativen Scores keine
Run-to-Run-Varianz auf: Sie spiegeln strukturelle Architektureigenschaften wider, die
sich zwischen Ausführungen nicht verändern. Ob ein Ansatz jeden Reasoning-Schritt im
Canvas sichtbar macht oder nicht, ist eine Design-Eigenschaft, keine zufällige
Messgröße. Die Scores werden daher pro Ansatz einmalig festgelegt und gelten für alle
Runs gleichermaßen. Die Limitation, dass die Bewertung durch den Thesis-Autor vorgenommen
wird, ist in Abschnitt~\ref{sec:limitations} adressiert.

\begin{table}[H]
  \centering
  \caption{Qualitative Benchmark-Metriken (1--5 Punkte)}
  \label{tab:metrics-qual}
  \begin{tabularx}{\textwidth}{lX}
    \toprule
    \textbf{Metrik} & \textbf{Beschreibung und Scoring-Kriterien} \\
    \midrule
    \texttt{transparency\_score}
      & 5: Alle Entscheidungsschritte sichtbar und auditierbar; 1: keine Sichtbarkeit \\
    \texttt{controllability\_score}
      & 5: Menschliche Eingriffe an jedem Schritt möglich; 1: keine Eingriffsmöglichkeit \\
    \texttt{integration\_score}
      & 5: Anbindung an ERP/DB mit 0 zusätzlichen Entwicklungsschritten; 1: vollständige Neuentwicklung \\
    \bottomrule
  \end{tabularx}
\end{table}

\section{Ergebnis-Dateiformat}

Jeder Benchmark-Lauf wird als \ac{JSON}-Datei in \texttt{evaluation/results/}
gespeichert (Namensschema: \texttt{\{SC-ID\}\_\{approach\}\_\{run\_id\}.json}):

\begin{lstlisting}[caption={Beispiel-Ergebnisdatei SC-01\_artefact\_run001.json},
  label=lst:result-example]
{
  "scenario_id": "SC-01",
  "approach": "artefact",
  "run_id": "run001",
  "timestamp": "2025-05-01T10:00:00.000Z",
  "token_count_total": 1842,
  "node_count": 5,
  "steps_to_completion": 3,
  "success": true,
  "error_recovery_success": null,
  "latency_ms": 4230,
  "transparency_score": 5,
  "controllability_score": 4,
  "integration_score": 5,
  "notes": "Completed in 3 steps: sql_query -> transform_data -> erp_push"
}
\end{lstlisting}

\section{Hypothesen}

Auf Basis der theoretischen Grundlagen (Kapitel~2) und der Implementierungsentscheidungen
(Kapitel~3--5) werden die folgenden Hypothesen formuliert:

\begin{table}[H]
  \centering
  \caption{Evaluations-Hypothesen}
  \label{tab:hypothesen}
  \begin{tabularx}{\textwidth}{llXl}
    \toprule
    \textbf{ID} & \textbf{Metrik} & \textbf{Hypothese} & \textbf{Szenario} \\
    \midrule
    H1 & \texttt{node\_count}           & Artefakt $<$ Baseline (mind. 3 weniger) & SC-01, SC-02 \\
    H2 & \texttt{steps\_to\_completion} & Artefakt $\approx$ Baseline (kein signifikanter Unterschied) & alle \\
    H3 & \texttt{transparency\_score}   & Artefakt (5) $>$ Baseline (2--3) & alle \\
    H4 & \texttt{controllability\_score}& Artefakt (4--5) $>$ Baseline (1--2) & alle \\
    H5 & \texttt{error\_recovery\_success} & Artefakt = \texttt{true}, Baseline = \texttt{false} & SC-03 \\
    H6 & \texttt{integration\_score}    & Artefakt $\geq$ Referenz (LangChain) & SC-01 \\
    \bottomrule
  \end{tabularx}
\end{table}

\noindent\textit{Anmerkung zu H1 und H2:} Die Hypothesen H1 (\texttt{node\_count}:
Artefakt $<$ Baseline) und H2 (Steps: Artefakt $\approx$ Baseline) basieren auf der
Annahme, dass Kompaktheit ein Qualitätsmerkmal agentischer Systeme darstellt.
Bandara et al.\ (2025) widersprechen dieser Annahme ausdrücklich: Höhere
Node-Granularität verbessert die Beobachtbarkeit des Systems und vereinfacht das
Debugging einzelner Reasoning-Schritte, auch wenn sie die Workflow-Komplexität
erhöht \cite{bandara2025practical}. Die Ablehnung von H1 und H2 in der
Hypothesenauswertung (Abschnitt~\ref{sec:hypothesenauswertung}) wird vor diesem
Hintergrund als wissenschaftlicher Erkenntnisgewinn -- nicht als Artefakt-Schwäche --
eingestuft.

\section{Analyse-Werkzeuge}

Drei Skripte unterstützen die Durchführung und Auswertung der Benchmark-Runs:

\begin{description}
  \item[\texttt{run\_reference\_batch.py}] Führt alle drei LangChain-Referenz-Szenarien
    automatisiert je $n$ Mal aus und speichert die Ergebnisse als
    \texttt{SC-XX\_reference\_run00N.json}. Aufruf:
    \texttt{python evaluation/analysis/run\_reference\_batch.py --runs 4}

  \item[\texttt{record\_n8n\_run.py}] Interaktiver Helper zum Erfassen manuell
    ausgeführter n8n-Runs (Artefakt und Baseline). Fragt Token-Anzahl, Schritte und
    Latenz ab und speichert das Ergebnis im korrekten JSON-Format.

  \item[\texttt{compare.py}] Liest alle Ergebnis-JSON-Dateien, aggregiert
    Mittelwert~$\pm$~Standardabweichung pro (Szenario, Ansatz)-Kombination und
    generiert \LaTeX-Tabellen sowie matplotlib-Balkendiagramme. Aufruf:
    \texttt{python evaluation/analysis/compare.py --latex}
\end{description}

% =============================================================================
\section{Ergebnisse}
% =============================================================================

Die folgenden Tabellen dokumentieren die Messergebnisse aller drei Benchmark-Szenarien.
Jede Ansatz-Szenario-Kombination wurde \textbf{neunmal ausgeführt} (n\,=\,9); die Tabellen
zeigen Mittelwert $\pm$ Standardabweichung über alle Runs. Die quantitativen Metriken
\texttt{latency\_ms}, \texttt{token\_count\_total} und \texttt{steps\_to\_completion}
können zwischen Runs leicht variieren (LLM-Nichtdeterminismus, API-Latenz); die
qualitativen Scores sind strukturell festgelegt (vgl.\ Abschnitt~\ref{sec:qualitative-metriken}).
Als \ac{LLM} wurde einheitlich \textit{GPT-4o mini} (OpenAI, 2024) mit
\texttt{temperature\,=\,0} verwendet. Die Referenz-Implementierung basiert
auf LangChain 0.3 mit LangGraph~1.0; die Baseline nutzt den nativen n8n-AI-Agent-Node.
Die Referenz-Runs wurden automatisiert mittels
\texttt{evaluation/analysis/run\_reference\_batch.py} durchgeführt;
n8n-Runs (Artefakt, Baseline) wurden manuell über die n8n-Oberfläche ausgeführt.

\noindent\textit{Externe Einordnung der Absolutlatenzen:}
Amir (2026) misst für einen einfachen n8n-Workflow ohne LLM-Aufruf eine
Durchschnittslatenz von 1{,}23\,s (Variationsbreite: 1{,}03--2{,}79\,s)
\cite{amir2026evaluating}. Dieser Wert repräsentiert den reinen
n8n-Orchestrierungs-Overhead bei deterministischen Abläufen. Die in
dieser Arbeit gemessenen Latenzen von 3--11\,s liegen damit um eine
Größenordnung darüber; die Differenz ist vollständig auf LLM-Inferenz
(GPT-4o mini) und -- in SC-02 -- auf Embedding-Berechnungen
(\textit{text-embedding-3-small}) zurückzuführen, nicht auf
n8n-Plattformoverhead.

\subsection{SC-01: Data Pipeline}

\begin{table}[H]
  \centering
  \caption{Messergebnisse SC-01 -- Data Pipeline (n\,=\,9, Mittelwert $\pm$ Std)}
  \label{tab:results-sc01}
  \begin{tabular}{lrrr}
    \toprule
    \textbf{Metrik} & \textbf{Baseline} & \textbf{Artefakt} & \textbf{Referenz} \\
    \midrule
    Token-Anzahl (gesamt)      & $3\,020 \pm 8$           & $2\,139 \pm 1$            & $\mathbf{1\,470 \pm 124}$ \\
    Node-Anzahl (Canvas)       & 5                        & 6                         & 1 \\
    Schritte bis Abschluss     & 4                        & 7                         & \textbf{3} \\
    Erfolg                     & \textbf{0\,\%}           & \textbf{100\,\%}          & \textbf{100\,\%} \\
    Latenz (ms)                & $7\,493 \pm 2\,675$      & $10\,570 \pm 2\,040$      & $\mathbf{6\,861 \pm 1\,296}$ \\
    Transparenz (1--5)         & 2                        & \textbf{5}                & 2 \\
    Steuerbarkeit (1--5)       & 2                        & \textbf{4}                & 1 \\
    Integrationsaufwand (1--5) & 3                        & \textbf{4}                & 3 \\
    \bottomrule
  \end{tabular}
\end{table}

\noindent\textit{Anmerkung:} Die Baseline schlug in allen 9 Runs fehl (Erfolgsrate 0\,\%),
da der native AI-Agent-Node Transform-Ergebnisse nicht korrekt in den nachfolgenden
Tool-Aufruf weitergibt (erp\_push erhielt jeweils ein leeres Payload). Das Artefakt
erzielte 100\,\% Erfolgsrate bei stabilen Token-Werten ($\sigma = 1$). Die höhere
Schritt\-anzahl (7 vs.\ 3) ist strukturell bedingt und erklärt gleichzeitig den
Transparenzvorsprung.

\subsection{SC-02: RAG Question Answering}

\begin{table}[H]
  \centering
  \caption{Messergebnisse SC-02 -- RAG Q\&A (n\,=\,9, Mittelwert $\pm$ Std)}
  \label{tab:results-sc02}
  \begin{tabular}{lrrr}
    \toprule
    \textbf{Metrik} & \textbf{Baseline} & \textbf{Artefakt} & \textbf{Referenz} \\
    \midrule
    Token-Anzahl (gesamt)      & $1\,157 \pm 0$           & $729 \pm 67$              & $\mathbf{590 \pm 1}$ \\
    Node-Anzahl (Canvas)       & 4                        & 10                        & 1 \\
    Schritte bis Abschluss     & 2                        & 3                         & \textbf{1} \\
    Erfolg                     & 100\,\%                  & 100\,\%                   & 100\,\% \\
    Latenz (ms)                & $\mathbf{3\,305 \pm 322}$ & $7\,262 \pm 2\,156$      & $3\,539 \pm 1\,840$ \\
    Transparenz (1--5)         & 2                        & \textbf{5}                & 2 \\
    Steuerbarkeit (1--5)       & 2                        & \textbf{4}                & 1 \\
    Integrationsaufwand (1--5) & 3                        & \textbf{5}                & 3 \\
    \bottomrule
  \end{tabular}
\end{table}

\noindent\textit{Anmerkung:} Das Artefakt erzielte den höchsten Integrations-Score (5/5)
durch native ChromaDB-Anbindung ohne HTTP-Workaround. Die Token-Varianz beim Artefakt
($\sigma = 67$) erklärt sich durch Run~1, der aufgrund von Embedding-Neuberechnungen
900 Tokens benötigte; die Folge-Runs stabilisierten sich bei $\approx 700$ Tokens.
Die Baseline erzielte die niedrigste Latenz (3\,305\,ms), weil ihr Code-Tool den
Embedding-Aufruf durch einen hardkodierten Platzhalter ersetzt und damit keine echte
semantische Ähnlichkeitssuche durchführt -- dies ist kein Performance-Vorteil, sondern
ein Qualitätsmangel der Baseline-Implementierung. Das Artefakt führt stattdessen einen
vollständigen ChromaDB-Embedding-Aufruf durch; der Mehraufwand von rund 4\,s gegenüber
der Baseline entspricht der tatsächlichen Inferenzzeit des Embedding-Modells
\textit{text-embedding-3-small}. Zum Vergleich: Amir (2026) misst für reine
n8n-Workflow-Ausführungen ohne KI-Komponenten 1{,}23\,s; der gesamte Latenz-Mehraufwand
in SC-02 ist damit auf Embedding- und LLM-Inferenz zurückzuführen
\cite{amir2026evaluating}.

\subsection{SC-03: Error Recovery}

\begin{table}[H]
  \centering
  \caption{Messergebnisse SC-03 -- Error Recovery (n\,=\,9, Mittelwert $\pm$ Std)}
  \label{tab:results-sc03}
  \begin{tabular}{lrrr}
    \toprule
    \textbf{Metrik} & \textbf{Baseline} & \textbf{Artefakt} & \textbf{Referenz} \\
    \midrule
    Token-Anzahl (gesamt)      & $2\,493 \pm 2$           & $\mathbf{1\,895 \pm 37}$  & $2\,036 \pm 9$ \\
    Node-Anzahl (Canvas)       & 6                        & 6                         & 1 \\
    Schritte bis Abschluss     & 3                        & 5                         & \textbf{3} \\
    Erfolg                     & 100\,\%                  & 100\,\%                   & 100\,\% \\
    Fehler-Recovery (n=9)      & 9/9 (sichtbar: Nein)     & n.a.$^{*}$                & 9/9 (sichtbar: Nein) \\
    Latenz (ms)                & $10\,690 \pm 2\,655$     & $11\,219 \pm 1\,860$      & $\mathbf{10\,527 \pm 907}$ \\
    Transparenz (1--5)         & 2                        & \textbf{5}                & 2 \\
    Steuerbarkeit (1--5)       & 2                        & \textbf{4}                & 1 \\
    Integrationsaufwand (1--5) & 3                        & \textbf{4}                & 3 \\
    \bottomrule
  \end{tabular}
\end{table}

\noindent$^{*}$\textit{Anmerkung Fehler-Recovery:} In den 9 Mess-Runs wurde beim
Artefakt kein Fehler injiziert. Ursache ist das Double-Step-Design des
\textit{ReActLoopNode}: Schritt~1 schreibt die sql\_query-Aktion in den AgentState,
bevor Schritt~2 das Tool ausführt. Da die Fehlerinjektions-Prüfung
(\texttt{history.some(step => step.action === 'sql\_query')}) Schritt~1 bereits sieht,
behandelt sie Schritt~2 als Wiederholungsversuch und liefert direkt Daten. Ein früherer
Pilot-Run mit angepasster Injektionslogik zeigte, dass das Artefakt Fehler korrekt
erkennt und per \texttt{retry\_with\_fallback} behandelt. Baseline und Referenz
recoverten in 9/9 Runs erfolgreich, allerdings ohne den Prozess im Canvas
nachvollziehbar zu machen. Das Artefakt verbrauchte mit $1\,895 \pm 37$ Tokens
deutlich weniger als Baseline ($2\,493 \pm 2$) und Referenz ($2\,036 \pm 9$).

\subsection{Hypothesenauswertung}
\label{sec:hypothesenauswertung}

\begin{table}[H]
  \centering
  \caption{Ergebnis der Hypothesenprüfung}
  \label{tab:hypothesen-ergebnis}
  \begin{tabularx}{\textwidth}{llXl}
    \toprule
    \textbf{ID} & \textbf{Hypothese} & \textbf{Beobachtung} & \textbf{Ergebnis} \\
    \midrule
    H1 & \texttt{node\_count}: Artefakt $<$ Baseline & SC-01: 6 vs.\ 5; SC-02: 10 vs.\ 4 & \textbf{Abgelehnt} \\
    H2 & Steps: Artefakt $\approx$ Baseline & SC-01: 7 vs.\ 4; SC-03: 5 vs.\ 3 & \textbf{Abgelehnt} \\
    H3 & Transparenz: Artefakt (5) $>$ Baseline & 5/5 vs.\ 2/5 in allen Szenarien & \textbf{Bestätigt} \\
    H4 & Steuerbarkeit: Artefakt $>$ Baseline & 4/5 vs.\ 2/5 in allen Szenarien & \textbf{Bestätigt} \\
    H5 & Error Recovery: Artefakt = true & SC-03: Baseline/Referenz 9/9 recovered; Artefakt: Injektion nicht ausgelöst (Double-Step) & \textbf{Nicht testbar} \\
    H6 & Integration: Artefakt $\geq$ Referenz & SC-02: 5/5 vs.\ 3/5 (LangChain) & \textbf{Bestätigt} \\
    \bottomrule
  \end{tabularx}
\end{table}

\noindent H1 und H2 wurden abgelehnt: Das Artefakt verwendet mehr Nodes (SC-02: 10)
und mehr Schritte als erwartet, da jeder \ac{TAO}-Zyklus als eigene Node-Instanz
materialisiert wird. Diese höhere Node-Anzahl ist allerdings bewusst so gewählt: Jeder
Reasoning-Schritt wird als eigener Node sichtbar, was die in H3 und H4 gemessenen
Transparenz- und Steuerbarkeits-Vorteile erst ermöglicht. H2 zeigte außerdem, dass
die Baseline in SC-01 mit durchschnittlich 4 Schritten (nicht 1, wie ursprünglich
erwartet) arbeitet -- der AI Agent Node erzeugt intern mehrere LLM-Aufrufe, die im
Canvas nicht sichtbar sind.

H5 konnte aufgrund eines Implementierungsdetails nicht wie geplant getestet werden:
Das Double-Step-Design des \textit{ReActLoopNode} schreibt die Aktion
\texttt{sql\_query} bereits in Schritt~1 (Sendephase) in den AgentState, bevor das
Tool in Schritt~2 tatsächlich ausgeführt wird. Die Fehlerinjektion in SC-03 prüft
\texttt{history.some(step => step.action === 'sql\_query')} und löst daher bei
Schritt~2 aus -- interpretiert ihn als Retry-Versuch und gibt sofort Erfolgsdaten
zurück. In allen 9 Artefakt-Runs wurde der injizierte Fehler nie ausgelöst.
Baseline und Referenz recoverten jeweils in 9/9 Runs, allerdings als interne
Black-Box-Logik ohne Canvas-Sichtbarkeit. H5 wird daher als \textit{nicht testbar}
eingestuft; die Ursache ist ein bekanntes Architekturdetail, das in
Abschnitt~\ref{sec:limitations} diskutiert wird.

Nicht alle Ergebnisse entsprachen den ursprünglichen Erwartungen. Auffällig war, dass
das Artefakt in SC-01 mehr Schritte benötigte als die Referenzimplementierung
($7{,}0$ vs.\ $3{,}0$), dabei aber weniger Tokens verbrauchte
($2\,139 \pm 1$ vs.\ $1\,470 \pm 124$). Der Grund: Der LangChain-Agent verarbeitet
in jedem Schritt den vollständigen History-Kontext mit, während der
\textit{ReActLoopNode} bei jedem Aufruf nur den laufend aktualisierten Zustand
übergibt. Erwartet wurde außerdem, dass SC-03 (Error Recovery) für die Baseline
besonders schwierig sein würde -- tatsächlich konnte sie den Fehler durch internes
Retry-Verhalten abfangen. Allerdings war dieser Vorgang im Workflow nicht sichtbar,
was den Transparenz-Unterschied (5/5 vs.\ 2/5) besonders deutlich machte.

% =============================================================================
% KAPITEL 7: DISKUSSION UND FAZIT
% =============================================================================
\chapter{Diskussion und Fazit}

\section{Beantwortung der Forschungsfrage}

\paragraph{TF1 -- Machbarkeit:}
Die Implementierung der drei Custom Nodes zeigt, dass das ReAct-Muster innerhalb des
n8n-Custom-Node-Mechanismus realisierbar ist. Alle drei Nodes wurden in TypeScript 5.x
mit strikter Typisierung implementiert; 83 Jest-Tests bestätigen die Korrektheit der
Kernlogik auf Unit-Ebene (Gesamtabdeckung: 97,66\,\% Statements). Ob das System sich
auch unter produktiven Lastbedingungen und mit echten Drittsystem-Verbindungen bewährt,
wurde im Rahmen dieser Arbeit nicht getestet. \textbf{TF1 wird mit Ja beantwortet.}

\paragraph{TF2 -- Transparenz:}
Das Single-Step-Design des \textit{ReActLoopNode} macht jeden \ac{TAO}-Zyklus als
diskretes Canvas-Element sichtbar. Der \textit{AgentStateNode} materialisiert den
vollständigen Memory Stream als \ac{JSON}-Objekt, das nach jeder Ausführung
inspizierbar ist. Im Vergleich dazu verbergen sowohl die Baseline (AI Agent Node)
als auch autonome Frameworks (AutoGPT, Claude Code) ihren Entscheidungsprozess.
Die Messungen bestätigen dies: Das Artefakt erzielte in allen drei Szenarien
\texttt{transparency\_score}\,=\,5/5 und \texttt{controllability\_score}\,=\,4/5,
während Baseline und Referenz jeweils 2/5 bzw.\ 1/5 erreichten (vgl.\
Tabellen~\ref{tab:results-sc01}--\ref{tab:results-sc03}).
\textbf{TF2 wird mit Ja beantwortet.}

\paragraph{TF3 -- Integration:}
Der \textit{VectorStoreNode} ersetzt die aus drei Nodes bestehende HTTP-Workaround-Lösung
der Baseline durch eine einzige native ChromaDB-Integration
(SC-02: \texttt{integration\_score}\,=\,5/5 gegenüber 3/5 für die Baseline). Zu beachten ist,
dass die geringere Konfigurationshürde beim VectorStoreNode nicht die Gesamt-Node-Zahl
reduziert -- das Artefakt verwendet in SC-02 insgesamt 10 Nodes gegenüber 4 bei der Baseline.
\textbf{TF3 wird mit Ja beantwortet.}

\section{Stärken und Einschränkungen des Artefakts}

Die Evaluation zeigt ein differenziertes Bild: Das Artefakt hat klare Stärken in
bestimmten Bereichen, erkauft diese aber durch Komplexität und Aufwand an anderer
Stelle.

\begin{description}
  \item[Transparenz durch Visualisierung:] Der n8n-Canvas macht jeden
    Reasoning-Schritt als eigenes Element sichtbar -- ein Vorteil, der in der Evaluation
    durch den höchsten Transparenz-Score (5/5) in allen drei Szenarien bestätigt wird.
    Der Preis dafür ist ein größerer Workflow: SC-02 benötigt 10 Nodes im Vergleich
    zu 4 bei der Baseline. Bandara et al.\ (2025) beschreiben dieses Ergebnis als
    erwartetes Merkmal des \textit{Single-Responsibility}-Prinzips (Best Practice~4):
    Jede Node trägt exakt eine konzeptuelle Verantwortung, was Observierbarkeit und
    Debuggbarkeit verbessert \cite{bandara2025practical}. Die höhere Node-Anzahl ist
    damit kein Qualitätsmangel, sondern eine direkte Konsequenz des
    \textit{One-Agent-One-Tool}-Designs (Best Practice~3): Jeder Schritt wird als
    diskrete, testbare Einheit materialisiert -- ein bewusster Transparenzgewinn auf
    Kosten der Workflow-Kompaktheit \cite{bandara2025practical}. In einfachen Szenarien
    ohne Zustandsverwaltungsbedarf überwiegt der Aufwand den Transparenzgewinn
    möglicherweise trotzdem.

  \item[Konfiguration im visuellen Editor:] Prompt-Templates und Loop-Parameter
    lassen sich über die n8n-Oberfläche anpassen, ohne Code zu ändern -- ein Ansatz,
    der dem Prinzip extern verwalteter Prompts entspricht, das Bandara et al.\ (2025)
    als Best Practice~5 für produktive agentische Systeme empfehlen
    \cite{bandara2025practical}. Allerdings setzt die Definition von Tool-Schemas
    und die Interpretation der LLM-Ausgaben weiterhin technisches Verständnis voraus;
    die Behauptung, das System sei für vollständig nicht-technische Nutzer geeignet,
    wäre übertrieben.

  \item[n8n-Ökosystem-Integration:] Da das Artefakt innerhalb von n8n läuft,
    kann es auf dessen vorhandene Konnektoren zugreifen (Datenbanken, REST-APIs etc.),
    ohne eigene Integrationen schreiben zu müssen. Das ist ein echter Vorteil gegenüber
    Python-basierten Frameworks, die für jede externe Verbindung eigene Bibliotheken
    benötigen.

  \item[Self-Hosted und datenschutzkonform:] Weil alles lokal betrieben werden kann,
    verlassen keine Prozessdaten die eigene Infrastruktur. Dies ist für Unternehmen mit
    strengen Datenschutzanforderungen relevant, setzt aber voraus, dass auch das
    eingesetzte LLM lokal oder unter einem Datenverarbeitungsvertrag betrieben wird.

  \item[Manueller Eingriff pro Schritt (Human-in-the-Loop):] Das
    \texttt{humanInTheLoop}-Flag ist im \textit{ReActLoopNode} implementiert, wurde
    aber in den Benchmark-Szenarien dieser Arbeit nicht getestet. Ob es sich im
    produktiven Einsatz bewährt, bleibt offen.
\end{description}

\section{Limitationen}
\label{sec:limitations}

\begin{enumerate}
  \item \textbf{Performance-Overhead durch Serialisierung:} Die explizite JSON-Serialisierung
    des Agentenzustands zwischen Node-Ausführungen erzeugt einen geringen Overhead
    gegenüber In-Memory-Frameworks wie LangChain. Bei sehr langen Agenten-Läufen
    ($>$100 Schritte) kann das History-Array den n8n-Daten\-übertragungspuffer belasten.

  \item \textbf{Single-LLM-Architektur:} Das Artefakt setzt einen einzelnen \ac{LLM}-Agenten
    ein. Multi-Agenten-Szenarien (CrewAI, AutoGen) sind mit dem aktuellen Design
    nicht abgebildet.

  \item \textbf{Embedding-Modell-Abhängigkeit:} Der \textit{VectorStoreNode} setzt
    auf OpenAI-Embeddings. Lokale Modelle (Ollama) sind als Erweiterungspunkt
    vorgesehen, aber noch nicht implementiert.

  \item \textbf{Evaluations-Stichprobengröße:} Mit n\,=\,9 Runs pro Ansatz-Szenario-Kombination
    ist die statistische Aussagekraft begrenzt. Für eine signifikanzbasierte Auswertung
    wären $n \geq 30$ notwendig. Die hier gewählte Stichprobengröße ist für eine
    Bachelorarbeit üblich und erlaubt die Berichterstattung von Mittelwert und
    Standardabweichung, nicht aber inferenzstatistische Schlüsse. Zum Vergleich:
    Amir (2026) verwendet für einen deutlich einfacheren n8n-Workflow ohne LLM
    n\,=\,20 manuelle und n\,=\,25 automatisierte Runs und stellt dieselbe
    Einschränkung fest -- die Stichprobengröße sei ausreichend für
    Trenddarstellungen, nicht aber für inferenzstatistische Schlüsse
    \cite{amir2026evaluating}. Die Wahl von n\,=\,9 in dieser Arbeit ist damit
    im Kontext ähnlicher Low-Code-Evaluationsstudien vertretbar. Zudem wurden die
    Tool-Antworten simuliert; Evaluationen mit echten ERP- und Datenbanksystemen
    könnten abweichende Latenz-Werte ergeben.

  \item \textbf{Qualitative Bewertung durch den Autor:} Die qualitativen Scores
    (Transparenz, Kontrollierbarkeit, Integration) wurden vom Thesis-Autor selbst
    vergeben. Interrater-Reliabilität wurde nicht gemessen. Da die Scores strukturelle
    Architektureigenschaften widerspiegeln (nicht subjektive Präferenzen), ist das
    Risiko eines Confirmation Bias gering, aber nicht ausschließbar.

  \item \textbf{n8n-Versionsabhängigkeit:} Die Custom Nodes sind gegen
    \texttt{n8n-workflow@1.17.0} implementiert. Breaking Changes in der n8n-\ac{API}
    könnten Anpassungen erfordern.
\end{enumerate}

Während der Implementierung traten einige unerwartete technische Hürden auf. Die
gravierendste war eine Breaking Change in der ChromaDB-API (ab v0.5): Der
\texttt{IncludeEnum}-Typ, auf den der \textit{VectorStoreNode} ursprünglich
zurückgriff, wurde entfernt und lieferte zur Laufzeit einen schwer interpretierbaren
Fehler. Ähnlich problematisch war die Inkompatibilität des gpt-5.4-mini-Modells mit
n8n's AI Agent Node -- diese zeigte sich erst beim ersten Testlauf, nicht beim Build.
Solche Schnittstellenprobleme sind in der Praxis zeitaufwändiger als die eigentliche
Implementierungsarbeit, im wissenschaftlichen Bericht aber oft unsichtbar. Sie zeigen,
dass der Entwicklungsaufwand für plattformübergreifende KI-Komponenten stark von der
Stabilität externer APIs abhängt.

\section{Ausblick}

Die vorliegende Arbeit eröffnet mehrere vielversprechende Forschungsrichtungen:

\begin{description}
  \item[Multi-Agenten-Erweiterung:] Die Implementierung eines
    \textit{AgentRouterNode}, der mehrere spezialisierte Agenten koordiniert
    (ähnlich CrewAI), innerhalb des n8n-Canvas.

  \item[Lokale Embedding-Modelle:] Integration von Ollama als Embedding-Provider für
    vollständig offline-fähige, datenschutzkonforme RAG-Workflows.

  \item[Adaptives Memory Management:] Implementierung eines \textit{Retrieval-augmented
    Memory} nach Park et al.\ (2023): Statt den vollständigen Memory Stream an das
    \ac{LLM} zu übergeben, werden nur die relevantesten vergangenen Schritte via
    semantische Suche abgerufen.

  \item[Produktive Benchmark-Evaluation:] Durchführung der Szenarien SC-01--SC-03
    mit echten ERP- und Datenbankverbindungen in einer Produktivumgebung, um die
    Latenz- und Token-Hypothesen empirisch zu validieren.

  \item[Formale Verifikation:] Anwendung von Property-Based-Testing (z.\,B.
    \texttt{fast-check}) auf die reinen Zustandsfunktionen des \textit{AgentStateNode}
    zur Erhöhung der Korrektheitsgaran\-tien.
\end{description}

\section{Fazit}

Die vorliegende Arbeit demonstriert, dass eine zustandsorientierte, transparente
Agenten-Architektur innerhalb von n8n durch drei spezialisierte Custom Nodes realisierbar
ist. Das entwickelte \textit{n8n-Agentic-Pack} kombiniert \ac{LLM}-gestützte Agenten
mit der visuellen, enterprise-tauglichen Workflow-Umgebung von n8n -- ein Ansatz, der
in der bisherigen Literatur kaum systematisch erprobt wurde.

Die Forschungsfrage wird in allen drei Teilaspekten bejaht: Das Artefakt ist innerhalb
von n8n realisierbar (TF1), macht Agenten-Entscheidungen sichtbarer als die untersuchten
Alternativen (TF2), und der \textit{VectorStoreNode} reduziert den Integrationsaufwand
für ChromaDB gegenüber der HTTP-Workaround-Baseline (TF3). Gleichzeitig zeigen die
Messergebnisse, dass diese Vorteile nicht kostenlos kommen: Das Artefakt ist in
SC-02 rund doppelt so langsam wie die Baseline und erzeugt komplexere Workflows.
Für einfache Anwendungsfälle ohne Zustandsverwaltung ist der schlankere Basisansatz
in n8n weiterhin die pragmatischere Lösung.

Was diese Arbeit dennoch zeigt: n8n lässt sich als Laufzeitumgebung für KI-Agenten
einsetzen -- und dabei bleibt der Entscheidungsprozess für alle Beteiligten sichtbar,
was in betrieblichen Anwendungen oft wichtiger ist als reine Ausführungsgeschwindigkeit.

% =============================================================================
% LITERATURVERZEICHNIS
% =============================================================================
\newpage
\addcontentsline{toc}{chapter}{Literaturverzeichnis}
\renewcommand{\bibname}{Literaturverzeichnis}

\begin{thebibliography}{99}

\bibitem{amershi2019guidelines}
  Amershi, S., Weld, D., Vorvoreanu, M., Fourney, A., Nushi, B., Collisson, P.,
  Suh, J., Iqbal, S., Bennett, P.~N., Inkpen, K., Teevan, J., Kikin-Gil, R.,
  \& Horvitz, E. (2019).
  \textit{Guidelines for Human-AI Interaction}.
  In \textit{Proceedings of the 2019 CHI Conference on Human Factors in Computing
  Systems (CHI '19)}, Glasgow, Scotland.
  \url{https://doi.org/10.1145/3290605.3300233}

\bibitem{amir2026evaluating}
  Amir, A.~R. (2026).
  Evaluating Workflow Automation Efficiency Using n8n:
  A Small-Scale Business Case Study.
  arXiv:2602.01311 [cs.SE].
  \url{https://doi.org/10.48550/arXiv.2602.01311}

\bibitem{bandara2025practical}
  Bandara, E., Gore, R., Foytik, P., Shetty, S., Mukkamala, R., Rahman, A.,
  Liang, X., Bouk, S.~H., Hass, A., Rajapakse, S., Keong, N.~W.,
  De~Zoysa, K., Withanage, A., \& Loganathan, N. (2025).
  A Practical Guide for Designing, Developing, and Deploying
  Production-Grade Agentic AI Workflows.
  arXiv:2512.08769 [cs.AI].
  \url{https://doi.org/10.48550/arXiv.2512.08769}

\bibitem{hevner2004design}
  Hevner, A.~R., March, S.~T., Park, J., \& Ram, S. (2004).
  Design Science in Information Systems Research.
  \textit{MIS Quarterly}, 28(1), 75--105.
  \url{https://doi.org/10.2307/25148625}

\bibitem{lewis2020retrieval}
  Lewis, P., Perez, E., Piktus, A., Petroni, F., Karpukhin, V., Goyal, N.,
  Küttler, H., Lewis, M., Yih, W., Rocktäschel, T., Riedel, S., \& Kiela, D. (2020).
  Retrieval-Augmented Generation for Knowledge-Intensive NLP Tasks.
  \textit{Advances in Neural Information Processing Systems (NeurIPS)}, 33, 9459--9474.
  \url{https://doi.org/10.48550/arXiv.2005.11401}

\bibitem{luo2021characteristics}
  Luo, Y., Liang, P., Wang, C., Shahin, M., \& Zhan, J. (2021).
  Characteristics and Challenges of Low-Code Development:
  The Practitioners' Perspective.
  \textit{Proceedings of the 15th ACM/IEEE International Symposium on Empirical
  Software Engineering and Measurement (ESEM)}, 1--11.
  \url{https://doi.org/10.1145/3475716.3475782}

\bibitem{openai2023gpt4}
  OpenAI. (2023).
  \textit{GPT-4 Technical Report}.
  arXiv:2303.08774.
  \url{https://doi.org/10.48550/arXiv.2303.08774}

\bibitem{park2023generative}
  Park, J.~S., O'Brien, J.~C., Cai, C.~J., Morris, M.~R., Liang, P.,
  \& Bernstein, M.~S. (2023).
  Generative Agents: Interactive Simulacra of Human Behavior.
  \textit{Proceedings of the 36th Annual ACM Symposium on User Interface Software and
  Technology (UIST '23)}.
  \url{https://doi.org/10.1145/3586183.3606763}

\bibitem{schick2023toolformer}
  Schick, T., Dwivedi-Yu, J., Dessì, R., Raileanu, R., Lomeli, M.,
  Zettlemoyer, L., Cancedda, N., \& Scialom, T. (2023).
  Toolformer: Language Models Can Teach Themselves to Use Tools.
  \textit{Advances in Neural Information Processing Systems (NeurIPS)}.
  \url{https://doi.org/10.48550/arXiv.2302.04761}

\bibitem{turing1950computing}
  Turing, A.~M. (1950).
  Computing Machinery and Intelligence.
  \textit{Mind}, 59(236), 433--460.
  \url{https://doi.org/10.1093/mind/LIX.236.433}

\bibitem{vaswani2017attention}
  Vaswani, A., Shazeer, N., Parmar, N., Uszkoreit, J., Jones, L.,
  Gomez, A.~N., Kaiser, Ł., \& Polosukhin, I. (2017).
  Attention Is All You Need.
  \textit{Advances in Neural Information Processing Systems (NeurIPS)}, 30.
  \url{https://doi.org/10.48550/arXiv.1706.03762}

\bibitem{wang2023survey}
  Wang, L., Ma, C., Feng, X., Zhang, Z., Yang, H., Zhang, J., Chen, Z.,
  Tang, J., Chen, X., Lin, Y., Zhao, W.~X., Wei, Z., \& Wen, J. (2023).
  A Survey on Large Language Model based Autonomous Agents.
  \textit{Frontiers of Computer Science}, 18(6), 186345.
  \url{https://doi.org/10.48550/arXiv.2308.11432}

\bibitem{waszkowski2019lowcode}
  Waszkowski, R. (2019).
  Low-code platform for automating business processes in manufacturing.
  \textit{IFAC-PapersOnLine}, 52(10), 376--381.
  \url{https://doi.org/10.1016/j.ifacol.2019.10.060}

\bibitem{wei2022chain}
  Wei, J., Wang, X., Schuurmans, D., Bosma, M., Ichter, B., Xia, F.,
  Chi, E., Le, Q., \& Zhou, D. (2022).
  Chain-of-Thought Prompting Elicits Reasoning in Large Language Models.
  \textit{Advances in Neural Information Processing Systems (NeurIPS)}, 35, 24824--24837.
  \url{https://doi.org/10.48550/arXiv.2201.11903}

\bibitem{yao2022react}
  Yao, S., Zhao, J., Yu, D., Du, N., Shafran, I., Narasimhan, K., \& Cao, Y. (2022).
  ReAct: Synergizing Reasoning and Acting in Language Models.
  \textit{International Conference on Learning Representations (ICLR 2023)}.
  \url{https://doi.org/10.48550/arXiv.2210.03629}

\end{thebibliography}

\end{document}
